% First Order Differential Equations With Separable Variables — Pure Mathematics with Mechanics Upper Sixth — Cours signé OnBuch+
% Official MINESEC syllabus. Compiler avec : tectonic PMM05-first-order-differential-equations-with-separable-variables.tex
\newcommand{\DOCMATIERE}{Pure Mathematics with Mechanics}
\newcommand{\DOCNIVEAU}{Upper Sixth}
\newcommand{\DOCMODULE}{Upper Sixth — Pure mathematics with mechanics}
\newcommand{\DOCLECON}{5}
\newcommand{\DOCTITRE}{First Order Differential Equations With Separable Variables}
% OnBuch+ — préambule « cours signés » ENRICHI (compilé avec tectonic / XeLaTeX)
% Système visuel « Pop » d'OnBuch+ : fond crème (jamais blanc), bordures encre
% épaisses, ombres « dures » décalées, titres Archivo Black en capitales.
% Version MATHÉMATIQUES : graphes (pgfplots/tikz), boîtes pédagogiques
% (définition, propriété, méthode, attention, le savais-tu, exercice résolu,
% exercice, TP, bilan), page de garde et signature OnBuch+ ancrée en bas.
%
% Macros attendues dans main.tex AVANT \input{preamble} :
%   \DOCMATIERE  \DOCNIVEAU  \DOCMODULE  \DOCLECON (numéro)  \DOCTITRE
\documentclass[11pt]{article}

\usepackage{fontspec}
\usepackage[english]{babel}
\usepackage[a4paper,margin=2cm,top=3.4cm,bottom=2.8cm,headheight=1.4cm,footskip=1.3cm]{geometry}
\usepackage{amsmath,amssymb,amsfonts}
\usepackage{mathtools} % \xmapsto, \coloneqq... souvent utilisés par les modèles
\usepackage{cancel} % simplifications barrées (bilans de forces)
% \abs{z} (module, valeur absolue) et \norm{u} (norme), utilisés par les modèles
\providecommand{\abs}{}\let\abs\relax\DeclarePairedDelimiter{\abs}{\lvert}{\rvert}
\providecommand{\norm}{}\let\norm\relax\DeclarePairedDelimiter{\norm}{\lVert}{\rVert}
\usepackage[table]{xcolor} % [table] : fournit aussi \rowcolors (alternance de lignes)
\usepackage{pagecolor}
\usepackage{tikz}
\usetikzlibrary{arrows.meta,positioning,calc,shapes.geometric,shapes.misc,shapes.arrows,decorations.pathmorphing,decorations.pathreplacing,decorations.markings,patterns,shadows,mindmap,trees,fit,backgrounds,matrix,intersections,angles,quotes}
\usepackage{pgfplots}
\usepackage[version=4]{mhchem} % équations nucléaires \ce{^{235}_{92}U}
\usepackage[european,siunitx]{circuitikz} % schémas électriques
\pgfplotsset{compat=1.18}
\usepgfplotslibrary{fillbetween,groupplots} % aires entre courbes (calcul intégral)
\usepackage{enumitem}
\usepackage{fancyhdr}
% [most] charge toutes les bibliothèques tcolorbox (listings, minted, raster,
% documentation...) et gonfle inutilement la mémoire TeX sur un document long
% (~300 boîtes) : on ne charge que ce qui est réellement utilisé.
\usepackage[skins,breakable]{tcolorbox}
\usepackage{graphicx}
\usepackage{colortbl}
\usepackage{booktabs}
\usepackage{tabularx}
\usepackage{diagbox} % case d'en-tête diagonale des tableaux à double entrée
% Colonnes extensibles centrée / gauche / droite (C, L, R), souvent utilisées par les modèles
\newcolumntype{C}{>{\centering\arraybackslash}X}
\newcolumntype{L}{>{\raggedright\arraybackslash}X}
\newcolumntype{R}{>{\raggedleft\arraybackslash}X}
\usepackage{array}
\usepackage{multicol}
\usepackage{siunitx}
% parse-numbers=false : accepte aussi \SI{2,5 \times 10^{-3}}{...}
\sisetup{locale=UK,output-decimal-marker={.},group-minimum-digits=4,parse-numbers=false}
\DeclareSIUnit\ohm{\ensuremath{\symup{\Omega}}} % U+2126 absent de la police
\DeclareSIUnit\division{div} % réglages d'oscilloscope (ms/div, V/div)
\DeclareSIUnit\tr{tr} % tours (vitesse de rotation, ex. \SI{1500}{\tr\per\minute})
\usepackage{qrcode}
% \degree : commande fréquemment utilisée par les modèles pour un angle ou une
% température en dehors de \SI{}{} (non fournie par les paquets chargés).
\providecommand{\degree}{^{\circ}}
% \qed (fin de démonstration), utilisé par les modèles sans amsthm
\providecommand{\qed}{\hfill$\square$}
% \barre : double barre des valeurs interdites dans un tableau de variations (tkz-tab)
\providecommand{\barre}{\ensuremath{\|}}
% environnement proof (amsthm non chargé) utilisé par les modèles
\ifdefined\proof\else\newenvironment{proof}[1][Proof]{\par\noindent\textit{#1.}\ }{\qed\par}\fi
\usepackage{lastpage}
\usepackage{hyperref}
\hypersetup{hidelinks}

% ============================================================
% Palette « Pop » OnBuch+ — mêmes hex que la classe P (plus_ui.dart)
% ============================================================
\definecolor{poporange}{HTML}{F59321}
\definecolor{popdark}{HTML}{E07A0C}
\definecolor{popcream}{HTML}{FFF6E8}
\definecolor{popink}{HTML}{1C1714}
\definecolor{poppurple}{HTML}{7A5AE0}
\definecolor{popgreen}{HTML}{2E9E6A}
\definecolor{poppink}{HTML}{E0457B}
\definecolor{popblue}{HTML}{2F7DE1}
\definecolor{popgold}{HTML}{C98A12}
\definecolor{popmuted}{HTML}{8A7F76}
\definecolor{popline}{HTML}{EADFD2}
\definecolor{popgray}{HTML}{8A7F76}
\colorlet{popgrey}{popgray}
% Alias tolérants (le modèle nomme parfois les couleurs différemment)
\colorlet{popwhite}{white}
\colorlet{popblack}{popink}
\colorlet{popred}{poppink}
\colorlet{popyellow}{popgold}
% Teintes claires pour fonds de tableaux / zones
\colorlet{popblueL}{popblue!10!white}
\colorlet{poporangeL}{poporange!14!white}
\colorlet{popgreenL}{popgreen!12!white}
\colorlet{poppurpleL}{poppurple!12!white}
\colorlet{poppinkL}{poppink!12!white}
\colorlet{popgoldL}{popgold!14!white}

\pagecolor{popcream}

% ============================================================
% Typographie
% ============================================================
\defaultfontfeatures{Ligatures=TeX}
\setmainfont{PlusJakartaSans-Medium.ttf}[Path=../../../fonts/,BoldFont=PlusJakartaSans-Bold.ttf]
% Maths en sans-serif (Fira Math) avec chiffres et lettres droites de Plus
% Jakarta, pour que les formules se fondent dans le texte.
\usepackage[math-style=ISO,bold-style=ISO]{unicode-math}
\setmathfont{FiraMath-Regular.otf}[Path=../../../fonts/]
\setmathfont{PlusJakartaSans-Medium.ttf}[Path=../../../fonts/,range={up/num,up/{Latin,latin},\mathup}]
% (icomma retiré : décimales avec point en anglais)
% Caractères mathématiques Unicode absents des polices de texte (Plus Jakarta,
% Archivo Black) : rendus par leur équivalent mathématique, en texte comme en
% formule — sinon ils s'affichent en carré vide (ex. « Arithmétique dans ℤ »).
\usepackage{newunicodechar}
\newunicodechar{ℝ}{\ensuremath{\mathbb{R}}}
\newunicodechar{ℤ}{\ensuremath{\mathbb{Z}}}
\newunicodechar{ℂ}{\ensuremath{\mathbb{C}}}
\newunicodechar{ℕ}{\ensuremath{\mathbb{N}}}
\newunicodechar{ℚ}{\ensuremath{\mathbb{Q}}}
\newunicodechar{𝔻}{\ensuremath{\mathbb{D}}}
\newunicodechar{α}{\ensuremath{\alpha}}
\newunicodechar{β}{\ensuremath{\beta}}
\newunicodechar{γ}{\ensuremath{\gamma}}
\newunicodechar{δ}{\ensuremath{\delta}}
\newunicodechar{ε}{\ensuremath{\varepsilon}}
\newunicodechar{θ}{\ensuremath{\theta}}
\newunicodechar{λ}{\ensuremath{\lambda}}
\newunicodechar{μ}{\ensuremath{\mu}}
\newunicodechar{σ}{\ensuremath{\sigma}}
\newunicodechar{τ}{\ensuremath{\tau}}
\newunicodechar{φ}{\ensuremath{\varphi}}
\newunicodechar{ψ}{\ensuremath{\psi}}
\newunicodechar{ω}{\ensuremath{\omega}}
\newunicodechar{Σ}{\ensuremath{\Sigma}}
% majuscules produites par \MakeUppercase dans les titres (α-aminés -> Α)
\newunicodechar{Α}{\ensuremath{\alpha}}
\newunicodechar{Β}{\ensuremath{\beta}}
\newunicodechar{Γ}{\ensuremath{\Gamma}}
\newunicodechar{Δ}{\ensuremath{\Delta}}
\newunicodechar{Ε}{\ensuremath{\varepsilon}}
\newunicodechar{Θ}{\ensuremath{\Theta}}
\newunicodechar{Λ}{\ensuremath{\Lambda}}
\newunicodechar{Μ}{\ensuremath{\mu}}
\newunicodechar{Τ}{\ensuremath{\tau}}
\newunicodechar{Φ}{\ensuremath{\Phi}}
\newunicodechar{Ψ}{\ensuremath{\Psi}}
\newunicodechar{Ω}{\ensuremath{\Omega}}
\newunicodechar{↦}{\ensuremath{\mapsto}}
\newunicodechar{∈}{\ensuremath{\in}}
\newunicodechar{∉}{\ensuremath{\notin}}
\newunicodechar{∧}{\ensuremath{\wedge}}
\newunicodechar{⇔}{\ensuremath{\Leftrightarrow}}
\newunicodechar{⟺}{\ensuremath{\Longleftrightarrow}}
\newunicodechar{⇒}{\ensuremath{\Rightarrow}}
\newunicodechar{⟹}{\ensuremath{\Longrightarrow}}
\newunicodechar{∘}{\ensuremath{\circ}}
\newunicodechar{∪}{\ensuremath{\cup}}
\newunicodechar{∩}{\ensuremath{\cap}}
\newunicodechar{≡}{\ensuremath{\equiv}}
\newunicodechar{⊂}{\ensuremath{\subset}}
\newunicodechar{⊥}{\ensuremath{\perp}}
\newunicodechar{∃}{\ensuremath{\exists}}
\newunicodechar{∀}{\ensuremath{\forall}}
\newunicodechar{∛}{\ensuremath{\sqrt[3]{\,}}}
\newunicodechar{∜}{\ensuremath{\sqrt[4]{\,}}}
\newunicodechar{⃗}{\ensuremath{{}^{\to}}}
\newunicodechar{ⁿ}{\ifmmode^{n}\else\textsuperscript{\ensuremath{n}}\fi}
\newunicodechar{ˣ}{\ifmmode^{x}\else\textsuperscript{\ensuremath{x}}\fi}
\newunicodechar{ᵇ}{\ifmmode^{b}\else\textsuperscript{\ensuremath{b}}\fi}
\newunicodechar{ʳ}{\ifmmode^{r}\else\textsuperscript{\ensuremath{r}}\fi}
\newunicodechar{ᵘ}{\ifmmode^{u}\else\textsuperscript{\ensuremath{u}}\fi}
\newunicodechar{ʸ}{\ifmmode^{y}\else\textsuperscript{\ensuremath{y}}\fi}
\newunicodechar{ᵃ}{\ifmmode^{a}\else\textsuperscript{\ensuremath{a}}\fi}
\newunicodechar{ᵉ}{\ifmmode^{e}\else\textsuperscript{\ensuremath{e}}\fi}
\newunicodechar{ᵏ}{\ifmmode^{k}\else\textsuperscript{\ensuremath{k}}\fi}
\newunicodechar{ᵖ}{\ifmmode^{p}\else\textsuperscript{\ensuremath{p}}\fi}
\newunicodechar{ᴺ}{\ifmmode^{N}\else\textsuperscript{\ensuremath{N}}\fi}
\newunicodechar{ᶜ}{\ifmmode^{c}\else\textsuperscript{\ensuremath{c}}\fi}
\newunicodechar{ʰ}{\ifmmode^{h}\else\textsuperscript{\ensuremath{h}}\fi}
\newunicodechar{ᵠ}{\ifmmode^{q}\else\textsuperscript{\ensuremath{q}}\fi}
\newunicodechar{ₙ}{\ifmmode_{n}\else\textsubscript{\ensuremath{n}}\fi}
\newunicodechar{ₐ}{\ifmmode_{a}\else\textsubscript{\ensuremath{a}}\fi}
\newunicodechar{ᵢ}{\ifmmode_{i}\else\textsubscript{\ensuremath{i}}\fi}
\newunicodechar{ᵣ}{\ifmmode_{r}\else\textsubscript{\ensuremath{r}}\fi}
\newunicodechar{ₖ}{\ifmmode_{k}\else\textsubscript{\ensuremath{k}}\fi}
\newunicodechar{ᵦ}{\ifmmode_{\beta}\else\textsubscript{\ensuremath{\beta}}\fi}
\newunicodechar{ₓ}{\ifmmode_{x}\else\textsubscript{\ensuremath{x}}\fi}
\newfontfamily\popdisplay{ArchivoBlack-Regular.ttf}[Path=../../../fonts/]
\newfontfamily\poplogo{SpaceGrotesk-Bold.ttf}[Path=../../../fonts/]
\newfontfamily\popsemi{PlusJakartaSans-SemiBold.ttf}[Path=../../../fonts/]
\newfontfamily\popxbold{PlusJakartaSans-ExtraBold.ttf}[Path=../../../fonts/]
\newfontfamily\popxboldls{PlusJakartaSans-ExtraBold.ttf}[Path=../../../fonts/,LetterSpace=3.0]

\setlist[enumerate]{leftmargin=1.7em,itemsep=5pt,topsep=4pt}
\setlist[itemize]{leftmargin=1.7em,itemsep=4pt,topsep=4pt,label={\color{poporange}\textbullet}}
\setlength{\parindent}{0pt}
\setlength{\parskip}{5pt}
\renewcommand{\arraystretch}{1.25}

% pgfplots : style Pop par défaut
\pgfplotsset{
  every axis/.append style={axis line style={popink,line width=1pt},tick style={popink},
    grid style={popline,line width=0.5pt},label style={font=\small},tick label style={font=\footnotesize}},
  popcurve/.style={poporange,line width=1.8pt,no markers,smooth},
}
% Même style disponible dans un \draw TikZ simple (hors pgfplots)
\tikzset{popcurve/.style={poporange,line width=1.8pt}}
\tikzset{popgrid/.style={popline,very thin}}
\colorlet{popcurve}{poporange} % le nom du style est parfois utilisé comme couleur

% ============================================================
% Pill / squiggle
% ============================================================
\newcommand{\poppill}[3]{%
  \tikz[baseline=-0.55ex]\node[draw=popink,line width=1.2pt,rounded corners=2.6pt,%
    fill=#2,inner xsep=6.5pt,inner ysep=3pt]{%
    \popxboldls\fontsize{7}{7}\selectfont\color{#3}\MakeUppercase{#1}};%
}
\newcommand{\popsquiggle}[1]{%
  \tikz[baseline=-0.4ex]{
    \draw[#1,line width=2.6pt,line cap=round]
      plot [smooth] coordinates {(0,0.09) (0.34,0.01) (0.68,0.21) (1.02,0.01) (1.36,0.10)};
  }%
}
% Mot-clé surligné (marqueur orange sous le texte)
\newcommand{\cle}[1]{{\popxbold\color{popdark}#1}}

% ============================================================
% En-tête / pied
% ============================================================
\pagestyle{fancy}
\fancyhf{}
\renewcommand{\headrulewidth}{0pt}
\renewcommand{\footrulewidth}{0pt}
\lhead{\raisebox{-0.15ex}{\poplogo\fontsize{13}{13}\selectfont\color{popink}OnBuch\color{poporange}+}}
\rhead{\poppill{\DOCMATIERE\ \textbullet\ \DOCNIVEAU\ \textbullet\ Chapter \DOCLECON}{white}{popink}}
\lfoot{\popxbold\fontsize{7.6}{7.6}\selectfont\color{popmuted}\MakeUppercase{Signed course OnBuch+}\ \textbullet\ {\popsemi\DOCTITRE}}
\rfoot{\tikz[baseline=-0.6ex]\node[rounded corners=8.5pt,fill=popink,minimum height=17pt,minimum width=17pt,inner xsep=5pt,inner sep=0]{\popxbold\fontsize{8}{8}\selectfont\color{popcream}\thepage\,/\,\pageref*{LastPage}};}

% ============================================================
% Titres
% ============================================================
\newcommand{\chaptitre}[2]{%
  \begin{tcolorbox}[enhanced,colback=poporange,colframe=popink,boxrule=2.6pt,arc=11pt,
      drop shadow=popink,left=15pt,right=15pt,top=12pt,bottom=11pt,before skip=3pt,after skip=18pt]
    \poppill{#2}{popink}{popcream}\par\vspace{9pt}
    {\raggedright\hyphenpenalty=10000\exhyphenpenalty=10000\popdisplay\fontsize{22}{24}\selectfont\color{popcream}\MakeUppercase{#1}\par}
  \end{tcolorbox}
}
% Section numérotée : pastille orange avec le numéro + titre Archivo
\newcounter{popsec}
\newcommand{\coursec}[1]{%
  \stepcounter{popsec}\setcounter{popsub}{0}%
  \par\vspace{16pt}\noindent
  \begin{minipage}{\linewidth}
  \tikz[baseline=-0.7ex]\node[draw=popink,line width=1.6pt,fill=poporange,rounded corners=4pt,minimum size=22pt,inner sep=0]{\popdisplay\fontsize{12}{12}\selectfont\color{popcream}\thepopsec};%
  \hspace{8pt}{\popdisplay\fontsize{15}{17}\selectfont\color{popink}\MakeUppercase{#1}}\par
  \vspace{4pt}\noindent{\color{poporange}\rule{36pt}{3.6pt}}
  \end{minipage}\par\nopagebreak\vspace{8pt}
}
\newcounter{popsub}
\newcommand{\courssub}[1]{%
  \stepcounter{popsub}%
  \par\vspace{9pt}\noindent{\popxbold\fontsize{11.5}{13}\selectfont\color{popink}{\color{poporange}\thepopsec.\thepopsub}\hspace{6pt}#1}\par\nopagebreak\vspace{4pt}
}

% ============================================================
% Boîtes pédagogiques — même recette : fond blanc, bordure épaisse colorée,
% ombre dure encre, badge plein. Titre optionnel [#1].
% ============================================================
\tcbset{popbase/.style={breakable,enhanced,colback=white,boxrule=2.3pt,arc=9pt,
  shadow={3pt}{-3pt}{0pt}{popink},left=14pt,right=14pt,top=11pt,bottom=11pt,before skip=11pt,after skip=15pt,
  fontupper=\popsemi\fontsize{10.6}{14.5}\selectfont\color{popink},
  fontlower=\popsemi\fontsize{10.6}{14.5}\selectfont\color{popink}}}
% Coche dessinée (le glyphe ✓ n'existe pas dans les polices du document)
\newcommand{\popcheck}[1]{\tikz[baseline=-0.5ex]\draw[#1,line width=1.6pt,line cap=round,line join=round](-0.13,0.0)--(-0.04,-0.09)--(0.13,0.1);}
\providecommand{\coche}{\popcheck{popgreen}}
\providecommand{\resultat}[1]{\par\smallskip\noindent\fcolorbox{popgreen}{popgreenL}{\parbox{\dimexpr\linewidth-2\fboxsep-2\fboxrule\relax}{\textbf{Result.} #1}}\par\smallskip}
\newcommand{\popboxhead}[3]{\poppill{#1}{#2}{white}\ifx\relax#3\relax\else\hspace{6pt}{\popxbold\fontsize{10.6}{12}\selectfont\color{popink}#3}\fi\par\vspace{7pt}}

\newtcolorbox{aretenir}[1][]{popbase,colframe=popgreen,before upper={\popboxhead{Key points}{popgreen}{#1}}}
\newtcolorbox{exemplebox}[1][]{popbase,colframe=poporange,before upper={\popboxhead{Exemple}{poporange}{#1}}}
\newtcolorbox{definition}[1][]{popbase,colframe=popblue,before upper={\popboxhead{Definition}{popblue}{#1}}}
\newtcolorbox{propriete}[1][]{popbase,colframe=poppurple,before upper={\popboxhead{Property}{poppurple}{#1}}}
\newtcolorbox{methode}[1][]{popbase,colframe=popgold,before upper={\popboxhead{Method}{popgold}{#1}}}
\newtcolorbox{attention}[1][]{popbase,colframe=poppink,before upper={\popboxhead{Warning}{poppink}{#1}}}
\newtcolorbox{experience}[1][]{popbase,colframe=popblue,colback=popblueL,before upper={\popboxhead{Experiment}{popblue}{#1}}}
% « Le savais-tu ? » : carte violette pleine (contexte, culture, Cameroun)
\newtcolorbox{savaistu}[1][]{popbase,colback=poppurple,colframe=popink,
  fontupper=\popsemi\fontsize{10.4}{14.2}\selectfont\color{white},
  before upper={\poppill{Did you know?}{white}{poppurple}\ifx\relax#1\relax\else\hspace{6pt}{\popxbold\color{white}#1}\fi\par\vspace{7pt}}}
% Exercice résolu : énoncé en haut, \tcblower puis la solution sur fond crème
\newcounter{popexo}\newcounter{popexr}
\newtcolorbox{exoresolu}[1][]{popbase,colframe=poporange,colbacklower=popcream,
  segmentation style={popink,line width=1.2pt,dashed},
  before upper={\stepcounter{popexr}\popboxhead{Worked example \thepopexr}{poporange}{#1}},
  before lower={\poppill{Solution}{popink}{popcream}\par\vspace{6pt}}}
% Exercice (non résolu en place) — la correction est dans la partie Corrigés
\newtcolorbox{exercice}[1][]{popbase,colframe=popink,
  before upper={\stepcounter{popexo}\popboxhead{Exercise \thepopexo}{popink}{#1}}}
% Corrigé
\newtcolorbox{corrige}[1][]{popbase,colframe=popgreen,colback=popgreenL,
  before upper={\popboxhead{Solution}{popgreen}{#1}}}
% Objectifs / prérequis / bilan
\newtcolorbox{objectifs}{popbase,colframe=popink,colback=poporangeL,
  before upper={\popboxhead{Chapter objectives}{popink}{}}}
\newtcolorbox{prerequis}{popbase,colframe=popink,colback=popblueL,
  before upper={\popboxhead{Prerequisites}{popblue}{}}}
\newtcolorbox{bilan}{popbase,colframe=popink,colback=white,boxrule=2.8pt,
  before upper={\popboxhead{Fiche bilan}{poporange}{L'essentiel à connaître pour l'examen}}}
% Figure encadrée avec légende
\newtcolorbox{popfigure}[1][]{enhanced,breakable=false,colback=white,colframe=popline,boxrule=1.4pt,arc=8pt,
  left=10pt,right=10pt,top=10pt,bottom=8pt,before skip=10pt,after skip=12pt,halign=center,
  lower separated=false,halign lower=center,
  fontlower=\popsemi\fontsize{9}{11.5}\selectfont\color{popmuted},
  #1}
\newcommand{\legende}[1]{\par\vspace{6pt}{\popsemi\fontsize{9}{11.5}\selectfont\color{popmuted}#1}}

% ============================================================
% Signature OnBuch+
% ============================================================
\newcommand{\poplogoinline}[1]{{\poplogo\fontsize{#1}{#1}\selectfont\color{popink}OnBuch\color{poporange}+}}
\newcommand{\signatureonbuch}{%
  \begin{tcolorbox}[enhanced,colback=popink,colframe=popink,boxrule=2.2pt,arc=10pt,
      drop shadow=poporange,left=14pt,right=14pt,top=10pt,bottom=10pt,before skip=0pt,after skip=0pt]
    \tikz[baseline=-0.7ex]\node[circle,fill=popgreen,draw=popcream,line width=1.3pt,minimum size=22pt,inner sep=0]{\popcheck{white}};%
    \hspace{9pt}\begin{minipage}[c]{0.62\linewidth}
      {\popxbold\fontsize{12}{13}\selectfont\color{popcream}Signed\ \textbullet\ {\poplogo OnBuch\color{poporange}+}}\par\vspace{2pt}
      {\raggedright\popsemi\fontsize{8}{10.5}\selectfont\color{popline}\MakeUppercase{OnBuch+ teaching team}\\\MakeUppercase{Follows the official MINESEC syllabus}\par}
    \end{minipage}\hfill
    {\poplogo\fontsize{22}{22}\selectfont\color{popcream}OnBuch\color{poporange}+}
  \end{tcolorbox}%
}

% ============================================================
% Page de garde
% #1 titre · #2 matière · #3 niveau · #4 module · #5 n° leçon · #6 durée ·
% #7 description / objectifs (texte ou itemize)
% ============================================================
\newcommand{\pagegarde}[7]{%
  \thispagestyle{empty}
  \newgeometry{a4paper,margin=1.7cm,top=1.6cm,bottom=1.4cm}%
  \begin{tikzpicture}[remember picture,overlay]
    \fill[poporange!18] ([xshift=-3.2cm,yshift=-3.6cm]current page.north east) circle (4.6cm);
    \fill[poppurple!14] ([xshift=2.4cm,yshift=7.6cm]current page.south west) circle (3.8cm);
  \end{tikzpicture}%
  {\poplogo\fontsize{30}{30}\selectfont\color{popink}OnBuch\color{poporange}+}\hfill
  \raisebox{0.6ex}{\poppill{Signed course \textbullet\ Premium}{popink}{popcream}}\par
  \vspace{1pt}\popsquiggle{poppurple}\par
  \vspace{8pt}
  \begin{tcolorbox}[enhanced,colback=poporange,colframe=popink,boxrule=2.8pt,arc=13pt,
      drop shadow=popink,left=18pt,right=18pt,top=12pt,bottom=13pt,after skip=10pt]
    \poppill{#2 \textbullet\ #3}{popink}{popcream}\hspace{5pt}\poppill{#4}{white}{popink}\par\vspace{8pt}
    {\popdisplay\fontsize{14}{14}\selectfont\color{popink}CHAPTER #5}\par\vspace{4pt}
    {\raggedright\hyphenpenalty=10000\exhyphenpenalty=10000\popdisplay\fontsize{25}{27}\selectfont\color{popcream}\MakeUppercase{#1}\par}
    \vspace{8pt}
    {\popxbold\fontsize{9.5}{9.5}\selectfont\color{popink}\MakeUppercase{Suggested time: #6}}
  \end{tcolorbox}
  \begin{tcolorbox}[enhanced,colback=white,colframe=popink,boxrule=2.2pt,arc=10pt,
      drop shadow=popink,left=16pt,right=16pt,top=10pt,bottom=10pt,after skip=8pt]
    \poppill{In this chapter}{poppurple}{white}\par\vspace{5pt}
    \popsemi\fontsize{9.3}{12.6}\selectfont\color{popink}#7
  \end{tcolorbox}
  \noindent\begin{minipage}[t]{0.487\linewidth}
    \begin{tcolorbox}[enhanced,colback=popgreen,colframe=popink,boxrule=2.2pt,arc=10pt,
        drop shadow=popink,left=11pt,right=11pt,top=10pt,bottom=10pt,height=3.1cm,valign=center]
      \tcbox[colback=white,colframe=popink,boxrule=1.3pt,arc=5pt,boxsep=0pt,nobeforeafter,
        left=3pt,right=3pt,top=3pt,bottom=3pt,valign=center]{\qrcode[height=1.9cm]{https://play.google.com/store/apps/details?id=com.luvvix.onbuch&hl=en}}%
      \hspace{8pt}\begin{minipage}[c]{0.5\linewidth}
        {\popdisplay\fontsize{11}{12.5}\selectfont\color{white}\MakeUppercase{Download OnBuch}}\par\vspace{4pt}
        {\raggedright\popsemi\fontsize{8.4}{11}\selectfont\color{white}Free \textbullet\ Google Play\\AI tutor, past papers, results\par}
      \end{minipage}
    \end{tcolorbox}
  \end{minipage}\hfill
  \begin{minipage}[t]{0.487\linewidth}
    \begin{tcolorbox}[enhanced,colback=poppurple,colframe=popink,boxrule=2.2pt,arc=10pt,
        drop shadow=popink,left=11pt,right=11pt,top=10pt,bottom=10pt,height=3.1cm,valign=center]
      \tikz[baseline=-0.7ex]\node[circle,fill=white,draw=popink,line width=1.3pt,minimum size=1.6cm,inner sep=0]{\popdisplay\fontsize{24}{24}\selectfont\color{poppurple}+};%
      \hspace{8pt}\begin{minipage}[c]{0.6\linewidth}
        {\popdisplay\fontsize{11}{12.5}\selectfont\color{white}\MakeUppercase{Go OnBuch+}}\par\vspace{4pt}
        {\raggedright\popsemi\fontsize{8.4}{11}\selectfont\color{white}All signed courses, unlimited AI tutor, PDF sheets — OnBuch+ tab in the app\par}
      \end{minipage}
    \end{tcolorbox}
  \end{minipage}
  \vspace{8pt}
  \popcontenu
  \vfill
  \signatureonbuch
  \restoregeometry
  \newpage
}

% Rangée « Ce cours contient » (page de garde)
\newcommand{\poptile}[3]{% #1 couleur #2 gros texte #3 légende
  \begin{tcolorbox}[enhanced,colback=#1,colframe=popink,boxrule=2pt,arc=9pt,shadow={2.5pt}{-2.5pt}{0pt}{popink},
      left=6pt,right=6pt,top=9pt,bottom=9pt,halign=center,valign=center,height=1.75cm,width=\linewidth,nobeforeafter]
    {\popdisplay\fontsize{12.5}{13}\selectfont\color{white}\MakeUppercase{#2}}\par\vspace{3pt}
    {\centering\popsemi\fontsize{7.8}{9.5}\selectfont\color{white}#3\par}
  \end{tcolorbox}}
\newcommand{\popcontenu}{%
  \par\noindent{\popxboldls\fontsize{7.5}{7.5}\selectfont\color{popmuted}THIS SIGNED COURSE CONTAINS}\par\vspace{7pt}
  \noindent\begin{minipage}[t]{0.235\linewidth}\poptile{popblue}{Course}{complete, illustrated, step by step}\end{minipage}\hfill
  \begin{minipage}[t]{0.235\linewidth}\poptile{poporange}{Methods}{and worked examples}\end{minipage}\hfill
  \begin{minipage}[t]{0.235\linewidth}\poptile{poppink}{Exercises}{exam style + solutions}\end{minipage}\hfill
  \begin{minipage}[t]{0.235\linewidth}\poptile{popgreen}{Summary sheet}{the essentials to remember}\end{minipage}\par}

% Sommaire
\newtcolorbox{sommaire}{enhanced,colback=white,colframe=popink,boxrule=2.2pt,arc=10pt,
  drop shadow=popink,left=18pt,right=18pt,top=13pt,bottom=13pt,before skip=6pt,after skip=20pt,
  before upper={\poppill{Contents}{poppurple}{white}\par\vspace{9pt}\popsemi\fontsize{10.6}{14.5}\selectfont\color{popink}}}

% Signature de fin de cours — poussée tout en bas de la dernière page
\newcommand{\findecours}{%
  \par\vfill\nopagebreak
  \begin{center}\popsquiggle{poporange}\end{center}\vspace{4pt}
  \signatureonbuch
}
\AtBeginDocument{\def\llbracket{\mathopen{[\mkern-3.5mu[}}\def\rrbracket{\mathclose{]\mkern-3.5mu]}}} % intervalles d'entiers (glyphe ⟦ absent de la police)
% Glyphes absents de Fira Math : redéfinis à partir de glyphes disponibles
\AtBeginDocument{%
  \def\setminus{\mathbin{\backslash}}\let\smallsetminus\setminus
  \def\vdots{\mathord{\vbox{\baselineskip3.2pt\lineskiplimit0pt\kern4pt\hbox{.}\hbox{.}\hbox{.}}}}%
  \def\ddots{\mathinner{\mkern1mu\raise6.5pt\hbox{.}\mkern2mu\raise3.6pt\hbox{.}\mkern2mu\raise0.7pt\hbox{.}\mkern1mu}}%
  \def\triangle{\mathord{\scalebox{0.9}{$\symup{\Delta}$}}}%
  \def\nabla{\mathord{\raisebox{\depth}{\scalebox{1}[-1]{$\symup{\Delta}$}}}}%
  \def\checkmark{\popcheck{popdark}}%
  \def\star{\mathbin{\ast}}\def\bigstar{\mathord{\ast}}%
  \def\diamond{\mathbin{\scalebox{0.6}{$\Diamond$}}}%
  \def\bigcup{\mathop{\mathchoice{\raisebox{-0.3em}{\scalebox{1.7}{$\cup$}}}{\scalebox{1.25}{$\cup$}}{\cup}{\cup}}\displaylimits}%
  \def\bigcap{\mathop{\mathchoice{\raisebox{-0.3em}{\scalebox{1.7}{$\cap$}}}{\scalebox{1.25}{$\cap$}}{\cap}{\cap}}\displaylimits}%
  \def\lesseqqgtr{\mathrel{\lessgtr}}\def\bigoplus{\mathop{\oplus}\displaylimits}%
}
\providecommand{\ssi}{if and only if} % abréviation fréquente
\AtBeginDocument{\providecommand{\wideparen}{\overparen}} % arc de cercle

%% FIGFIX-BEGIN v5 — correctifs globaux des figures (docs/onbuch-generation/tools/figfix)
\usepackage{adjustbox}\usepackage{etoolbox}\tcbuselibrary{raster}
% toute figure TikZ/pgfplots plus large que la ligne est réduite à la largeur disponible
\BeforeBeginEnvironment{tikzpicture}{\begin{adjustbox}{max width=\linewidth}}
\AfterEndEnvironment{tikzpicture}{\end{adjustbox}}
% légendes pgfplots lisibles par-dessus les courbes
\pgfplotsset{every axis/.append style={legend style={fill=white,fill opacity=0.92,text opacity=1,draw=black!15,inner sep=3pt}}}
% étiquettes d'arêtes (syntaxe edge["..."]) avec fond blanc
\tikzset{every edge quotes/.append style={fill=white,inner sep=1.2pt}}
%% FIGFIX-END
\begin{document}

\pagegarde{First Order Differential Equations With Separable Variables}{Pure Mathematics with Mechanics}{Upper Sixth}{Upper Sixth — Pure mathematics with mechanics}{5}{6 periods}{This chapter introduces first‑order differential equations with separable variables, guiding students from the basic definition through recognition, separation, integration, and the use of initial conditions to obtain particular solutions. It concludes with concrete applications drawn from physics, biology and everyday Cameroonian contexts.
\vspace{4pt}
\begin{multicols}{2}\small
\begin{itemize}
\item Define a first‑order differential equation and distinguish between general and particular solutions.
\item Recognise whether a given first‑order equation is separable.
\item Apply the separation‑of‑variables technique correctly, including handling constant solutions.
\item Integrate both sides of a separated equation using appropriate integration methods.
\item Express the general solution in implicit or explicit form and verify it by differentiation.
\item Use an initial condition to determine the arbitrary constant and write the particular solution.
\end{itemize}\end{multicols}}

\begin{sommaire}
\begin{multicols}{2}
\begin{enumerate}
\item Meaning of First Order Differential Equations
\item Recognition of Separable Equations
\item Separation of Variables
\item Integration to Obtain General Solutions
\item Particular Solutions and Applications
\item Integration activity
\item Methods and worked examples
\item Exercises
\item Solutions to the exercises
\item Summary sheet
\end{enumerate}
\end{multicols}
\end{sommaire}

\begin{objectifs}
\begin{itemize}
\item Define a first‑order differential equation and distinguish between general and particular solutions.
\item Recognise whether a given first‑order equation is separable.
\item Apply the separation‑of‑variables technique correctly, including handling constant solutions.
\item Integrate both sides of a separated equation using appropriate integration methods.
\item Express the general solution in implicit or explicit form and verify it by differentiation.
\item Use an initial condition to determine the arbitrary constant and write the particular solution.
\item Model real‑world phenomena (cooling, population growth, mixing) with separable differential equations.
\item Interpret the mathematical solution in the context of the original problem and discuss its domain of validity.
\end{itemize}
\end{objectifs}

\begin{prerequis}
\begin{itemize}
\item Differentiation rules (power, exponential, trigonometric, chain rule).
\item Integration techniques: substitution, partial fractions, standard integrals.
\item Algebraic manipulation of equations and handling of absolute values.
\item Basic understanding of functions, graphs, and slope fields.
\item Familiarity with the concept of a constant of integration.
\end{itemize}
\end{prerequis}

\newpage

\coursec{Meaning of First Order Differential Equations}

Imagine a steaming pot of ndolé left on the kitchen table in Yaoundé. As it cools, its temperature drops faster at first and then more slowly, following a natural law that can be described by a first‑order differential equation. In the same way, the spread of a rumour in a village or the growth of a cassava crop can be modelled with equations where the rate of change depends on the current state. This chapter equips you to translate such real‑life situations into mathematics and to solve them.

\courssub{What is a Differential Equation?}

A \cle{differential equation} is an equation that relates an unknown function to one or more of its derivatives. The function usually represents a quantity that varies (temperature, population, velocity, …) and the derivatives represent its rates of change.

\begin{definition}[Differential Equation]
A \textbf{differential equation} is an equation involving an unknown function \(y = y(x)\) and its derivatives \(\frac{dy}{dx}, \frac{d^2y}{dx^2}, \dots\). The \textbf{order} of the differential equation is the order of the highest derivative that appears. The \textbf{degree} is the power of the highest derivative after the equation has been made rational and polynomial in all derivatives.
\end{definition}

For example:
\[
\frac{dy}{dx} = 3x^2 \qquad\text{(order 1, degree 1)}
\]
\[
\left(\frac{d^2y}{dx^2}\right)^2 + y = \sin x \qquad\text{(order 2, degree 2)}
\]

\courssub{Order and Degree}

The \textbf{order} tells us how many times we must differentiate the unknown function to obtain the equation. The \textbf{degree} is only defined when the equation is a polynomial in the derivatives; it is the exponent of the highest‑order derivative.

\begin{attention}[Order vs. Degree]
A very common mistake is to confuse the order with the degree. Remember:
\begin{itemize}
\item \textbf{Order} = highest derivative present (first, second, …).
\item \textbf{Degree} = power of that highest derivative \emph{after} clearing fractions and radicals.
\end{itemize}
For instance, \(\displaystyle \left(\frac{dy}{dx}\right)^3 + y = x\) has order 1 and degree 3, whereas \(\displaystyle \frac{d^2y}{dx^2} + \sqrt{\frac{dy}{dx}} = 0\) has order 2 but its degree is not defined (because of the square root).
\end{attention}

\courssub{First‑Order Differential Equations}

A \cle{first‑order differential equation} involves only the first derivative \(\frac{dy}{dx}\). It can always be written in the form
\[
\frac{dy}{dx} = f(x, y)
\]
where \(f\) is a given function of the independent variable \(x\) and the dependent variable \(y\). The variable \(x\) is the \textbf{independent variable} (often time or position) and \(y\) is the \textbf{dependent variable} (the quantity we want to find).

Simple examples:
\begin{itemize}
\item \(\displaystyle \frac{dy}{dx} = 3x^2\) — the rate of change depends only on \(x\).
\item \(\displaystyle \frac{dy}{dx} = y\) — the rate of change is proportional to the current value (exponential growth/decay).
\item \(\displaystyle \frac{dy}{dx} = x + y\) — the rate of change depends on both \(x\) and \(y\).
\end{itemize}

\courssub{General and Particular Solutions}

When we solve a first‑order differential equation, we obtain a \cle{general solution} — a family of functions containing one arbitrary constant (usually denoted \(C\)). Each choice of \(C\) gives a \cle{particular solution}, which is a single curve in the \(xy\)-plane.

Geometrically, the differential equation \(\frac{dy}{dx} = f(x, y)\) assigns a slope \(f(x, y)\) to every point \((x, y)\) in the plane. The general solution is the family of all curves whose tangent at each point has exactly that slope. A particular solution is the unique curve from that family that passes through a given point \((x_0, y_0)\) — this point is called an \cle{initial condition}.

\courssub{Slope Fields (Direction Fields)}

A \cle{slope field} (or \cle{direction field}) is a graphical tool that helps us visualise the behaviour of solutions without solving the equation analytically. At a grid of points \((x_i, y_j)\) we draw a short line segment whose slope equals \(f(x_i, y_j)\). The collection of these segments shows the “flow” of the solutions.

\begin{methode}[Constructing a Slope Field]
\begin{enumerate}
\item Choose a rectangular grid of points \((x_i, y_j)\) in the region of interest.
\item At each grid point compute the slope \(m_{ij} = f(x_i, y_j)\).
\item Draw a short line segment centred at \((x_i, y_j)\) with slope \(m_{ij}\) (length can be constant or scaled).
\item Optionally, sketch a few integral curves that follow the direction of the segments.
\end{enumerate}
\end{methode}

\begin{popfigure}
\begin{tikzpicture}
\begin{axis}[
    width=12cm,
    height=8cm,
    axis lines=middle,
    xlabel=$x$, ylabel=$y$,
    xmin=-3, xmax=3,
    ymin=-3, ymax=3,
    xtick={-3,-2,-1,0,1,2,3},
    ytick={-3,-2,-1,0,1,2,3},
    grid=major,
    grid style={gray!30},
    view={0}{90},
]
\addplot3[
    quiver={
        u={1},
        v={x - y},
        scale arrows=0.3,
    },
    -stealth,
    samples=15,
    samples y=15,
    domain=-3:3,
    domain y=-3:3,
    color=popblue,
] {0};
\end{axis}
\end{tikzpicture}
\legende{Slope field for \(\frac{dy}{dx} = x - y\). The short arrows indicate the direction of the solution curves at each point.}
\end{popfigure}

The slope field above for \(\frac{dy}{dx} = x - y\) shows that solutions tend to approach the line \(y = x\) as \(x\) increases. This qualitative insight is often valuable before attempting an analytic solution.

\courssub{Integral Curves}

The curves that are everywhere tangent to the slope field are called \cle{integral curves}. For the simple equation \(\frac{dy}{dx} = y\) the general solution is \(y = Ce^x\). Plotting several members of this family on the same axes illustrates how the arbitrary constant \(C\) shifts the curve vertically.

\begin{popfigure}
\begin{tikzpicture}
\begin{axis}[
    width=12cm,
    height=8cm,
    axis lines=middle,
    xlabel=$x$, ylabel=$y$,
    xmin=-2, xmax=2,
    ymin=-5, ymax=5,
    xtick={-2,-1,0,1,2},
    ytick={-5,-4,-3,-2,-1,0,1,2,3,4,5},
    grid=major,
    grid style={gray!30},
    legend style={at={(0.02,0.98)}, anchor=north west, font=\small},
]
\addplot[domain=-2:2, samples=100, popblue, thick] {exp(x)};
\addlegendentry{$C=1$}
\addplot[domain=-2:2, samples=100, poporange, thick] {2*exp(x)};
\addlegendentry{$C=2$}
\addplot[domain=-2:2, samples=100, popgreen, thick] {0.5*exp(x)};
\addlegendentry{$C=0.5$}
\addplot[domain=-2:2, samples=100, poppurple, thick] {-exp(x)};
\addlegendentry{$C=-1$}
\addplot[domain=-2:2, samples=100, poppink, thick] {-2*exp(x)};
\addlegendentry{$C=-2$}
\end{axis}
\end{tikzpicture}
\legende{Integral curves of \(\frac{dy}{dx} = y\) for several values of the constant \(C\). All curves have the same exponential shape, differing only by a vertical scaling factor.}
\end{popfigure}

\courssub{Initial Conditions and Particular Solutions}

An \cle{initial condition} specifies the value of the unknown function at a particular point: \(y(x_0) = y_0\). Substituting this condition into the general solution determines the constant \(C\) uniquely, yielding a particular solution.

\begin{propriete}[Existence and Uniqueness (Brief Statement)]
If the function \(f(x, y)\) and its partial derivative \(\frac{\partial f}{\partial y}\) are continuous in a rectangle containing the point \((x_0, y_0)\), then the initial value problem
\[
\frac{dy}{dx} = f(x, y), \qquad y(x_0) = y_0
\]
has a unique solution defined on some interval around \(x_0\).
\end{propriete}

This theorem guarantees that, under mild conditions, the physical situations we model (cooling ndolé, spreading rumour, growing cassava) have a well‑defined future evolution once the present state is known.

\courssub{Worked Example}

\begin{exoresolu}[Finding a Particular Solution]
\textbf{Statement:} Solve the initial value problem
\[
\frac{dy}{dx} = 3x^2, \qquad y(1) = 4.
\]

\tcblower
\textbf{Detailed solution:}
\begin{enumerate}
\item \textbf{Recognise the equation.} It is a first‑order differential equation where the right‑hand side depends only on \(x\).
\item \textbf{Integrate both sides with respect to \(x\).}
\[
y = \int 3x^2 \, dx = x^3 + C,
\]
where \(C\) is an arbitrary constant. This is the \emph{general solution}.
\item \textbf{Apply the initial condition \(y(1)=4\).}
\[
4 = 1^3 + C \quad\Longrightarrow\quad C = 3.
\]
\item \textbf{Write the particular solution.}
\[
y = x^3 + 3.
\]
\item \textbf{Check.} \(\frac{dy}{dx} = 3x^2\) and \(y(1)=1+3=4\). Both the differential equation and the initial condition are satisfied.
\end{enumerate}
\end{exoresolu}

\courssub{Common Mistakes}

\begin{attention}[Frequent Errors]
\begin{itemize}
\item \textbf{Forgetting the constant of integration.} Every indefinite integration introduces an arbitrary constant; omitting it loses the general solution.
\item \textbf{Confusing the independent and dependent variables.} In \(\frac{dy}{dx}=f(x,y)\), \(x\) is independent, \(y\) is dependent. Swapping them changes the problem entirely.
\item \textbf{Misinterpreting the slope field.} The short segments represent the \emph{direction} of the solution curve at that point, not the curve itself. Do not join the tips of the segments; instead, sketch curves that are everywhere tangent to them.
\item \textbf{Assuming a unique solution without checking continuity.} The existence‑uniqueness theorem requires \(f\) and \(\frac{\partial f}{\partial y}\) to be continuous near the initial point. If they are not, multiple solutions or no solution may exist.
\end{itemize}
\end{attention}

\courssub{Summary of Key Points}

\begin{aretenir}[Key Points for This Section]
\begin{itemize}
\item A \textbf{differential equation} relates an unknown function to its derivatives.
\item The \textbf{order} is the highest derivative; the \textbf{degree} is the power of that derivative (when defined).
\item A \textbf{first‑order} equation has the form \(\frac{dy}{dx}=f(x,y)\).
\item The \textbf{general solution} contains one arbitrary constant and represents a family of curves (integral curves).
\item A \textbf{particular solution} is obtained by imposing an \textbf{initial condition} \(y(x_0)=y_0\).
\item \textbf{Slope fields} give a visual picture of the direction of solutions at each point.
\item Under continuity of \(f\) and \(\partial f/\partial y\), the initial value problem has a \textbf{unique} solution.
\end{itemize}
\end{aretenir}

\coursec{Recognition of Separable Equations}

In the previous section we learned what a first‑order differential equation is and how it arises from modelling rates of change.  We now turn to a particularly friendly family of such equations: those whose variables can be \cle{separated}.  Recognising this structure is the first decisive step towards an explicit solution.

\courssub{What Makes a Differential Equation Separable?}

\begin{definition}[Separable First‑Order Differential Equation]
A first‑order differential equation is called \textbf{separable} if it can be written in the form
\[
\frac{dy}{dx}=g(x)\,h(y)
\]
or, equivalently,
\[
M(x)\,N(y)\,dx + P(x)\,Q(y)\,dy = 0,
\]
where the functions of \(x\) and the functions of \(y\) appear only as products.  After a simple algebraic rearrangement we obtain
\[
\frac{1}{h(y)}\,dy = g(x)\,dx \qquad\text{or}\qquad f(x)\,dx = g(y)\,dy.
\]
\end{definition}

The key idea: \emph{all terms containing \(x\) (and \(dx\)) can be gathered on one side, and all terms containing \(y\) (and \(dy\)) on the other}.  Once this is achieved, each side can be integrated independently.

\courssub{Algebraic Test for Separability}

Given a differential equation, follow these algebraic steps to decide whether it is separable:

\begin{enumerate}
\item Write the equation explicitly as \(\dfrac{dy}{dx} = F(x,y)\).
\item Attempt to factor \(F(x,y)\) into a product \(g(x)\,h(y)\).  This may require elementary algebraic manipulations: expanding brackets, cancelling common factors, or rewriting a quotient as a product.
\item If such a factorisation exists, the equation is separable.  Divide both sides by \(h(y)\) (provided \(h(y)\neq 0\)) and multiply by \(dx\) to obtain the separated form
\[
\frac{1}{h(y)}\,dy = g(x)\,dx.
\]
\item If no factorisation into a pure product of a function of \(x\) and a function of \(y\) is possible, the equation is \emph{not} separable (though it might become separable after a suitable substitution — see the tip at the end of this section).
\end{enumerate}

\courssub{Examples of Separable Equations}

\begin{exemplebox}[Separable: \(dy/dx = x\,y\)]
\[
\frac{dy}{dx}=x\,y \quad\Longrightarrow\quad \frac{1}{y}\,dy = x\,dx.
\]
Here \(g(x)=x\) and \(h(y)=y\).  The variables are already separated after dividing by \(y\).
\end{exemplebox}

\begin{exemplebox}[Separable: \(dy/dx = \dfrac{\sin x}{y^2}\)]
\[
\frac{dy}{dx}=\frac{\sin x}{y^2} \quad\Longrightarrow\quad y^2\,dy = \sin x\,dx.
\]
We have \(g(x)=\sin x\) and \(h(y)=1/y^2\).  Multiplying by \(y^2\,dx\) yields the separated form.
\end{exemplebox}

\begin{exemplebox}[Separable after simplification: \(\dfrac{dy}{dx} = \dfrac{x^2y + y}{x}\)]
\[
\frac{dy}{dx} = \frac{y(x^2+1)}{x} = \underbrace{\frac{x^2+1}{x}}_{g(x)}\;\underbrace{y}_{h(y)}.
\]
Factorising the numerator reveals the product structure.
\end{exemplebox}

\courssub{Examples of Non‑Separable Equations}

\begin{exemplebox}[Not separable: \(dy/dx = x + y\)]
The right‑hand side is a \emph{sum} of a function of \(x\) and a function of \(y\).  It cannot be expressed as a product \(g(x)h(y)\) (except in trivial cases where one factor is constant, which would make the equation either purely \(x\)‑dependent or purely \(y\)‑dependent).
\end{exemplebox}

\begin{exemplebox}[Not separable: \(dy/dx = \dfrac{xy}{x+y}\)]
The denominator \(x+y\) mixes the variables in a way that prevents factorisation into a product of a function of \(x\) alone and a function of \(y\) alone.
\end{exemplebox}

\courssub{Hidden Separability — Simple Algebraic Manipulation}

Sometimes an equation does not look separable at first glance, but a little algebra reveals the structure.

\begin{exemplebox}[Hidden separability: \(\dfrac{dy}{dx} = \dfrac{y^2-1}{x^2-1}\)]
\[
\frac{dy}{dx} = \frac{y^2-1}{x^2-1}
= \underbrace{\frac{1}{x^2-1}}_{g(x)}\;\underbrace{(y^2-1)}_{h(y)}.
\]
The quotient of two differences of squares is already a product of a function of \(x\) and a function of \(y\).
\end{exemplebox}

\begin{exemplebox}[Hidden separability: \(\dfrac{dy}{dx} = \dfrac{e^{x+y}}{x}\)]
\[
\frac{dy}{dx} = \frac{e^x e^y}{x} = \underbrace{\frac{e^x}{x}}_{g(x)}\;\underbrace{e^y}_{h(y)}.
\]
Using the law of exponents \(e^{x+y}=e^x e^y\) exposes the separable form.
\end{exemplebox}

\courssub{Flowchart for Testing Separability}

\begin{popfigure}
\begin{tikzpicture}[
  node distance=1.8cm and 2.5cm,
  startstop/.style={rectangle, rounded corners, minimum width=3cm, minimum height=1cm, text centered, draw=popblue, fill=popblueL, thick},
  process/.style={rectangle, minimum width=4cm, minimum height=1cm, text centered, draw=popdark, fill=white, thick},
  decision/.style={diamond, minimum width=4cm, minimum height=1.2cm, text centered, draw=poporange, fill=poporange!20, thick, aspect=2},
  arrow/.style={->, thick, popdark},
  yesno/.style={midway, fill=white, inner sep=2pt, font=\small}
]
\node (start) [startstop] {Start: $\frac{dy}{dx}=F(x,y)$};
\node (factor) [process, below of=start] {Try to factor $F(x,y)$ as $g(x)h(y)$};
\node (decide) [decision, below of=factor] {Success?};
\node (separable) [process, below left of=decide, xshift=-2cm, yshift=-0.5cm, align=center]{Equation is separable\\Write $\frac{1}{h(y)}dy=g(x)dx$};
\node (notsep) [process, below right of=decide, xshift=2cm, yshift=-0.5cm, align=center]{Not separable\\Consider substitution or other methods};
\node (end) [startstop, below of=decide, yshift=-3.5cm] {Proceed to integration};

\draw [arrow] (start) -- (factor);
\draw [arrow] (factor) -- (decide);
\draw [arrow] (decide) -- node[yesno] {Yes} (separable);
\draw [arrow] (decide) -- node[yesno] {No} (notsep);
\draw [arrow] (separable) |- (end);
\draw [arrow] (notsep) |- (end);
\end{tikzpicture}
\legende{Flowchart: steps to test whether a first‑order DE is separable.}
\end{popfigure}

\courssub{Method: Step‑by‑Step Recognition}

\begin{methode}[Recognising a Separable Differential Equation]
\begin{enumerate}
\item Write the DE in the explicit form \(\dfrac{dy}{dx}=F(x,y)\).
\item Inspect \(F(x,y)\): can it be written as a product \(g(x)\,h(y)\)?  Use algebraic identities (factorisation, exponent laws, trigonometric identities) if necessary.
\item If yes, the DE is separable.  Divide both sides by \(h(y)\) (noting any values where \(h(y)=0\)) and multiply by \(dx\) to obtain \(\dfrac{1}{h(y)}\,dy = g(x)\,dx\).
\item If no, the DE is not separable.  Record this and consider whether a substitution (e.g. \(v=y/x\) for homogeneous equations) might render it separable.
\end{enumerate}
\end{methode}

\courssub{Important Warning: Lost Constant Solutions}

\begin{attention}[Dividing by \(h(y)\) May Lose Solutions]
When we divide by \(h(y)\) we implicitly assume \(h(y)\neq 0\).  If \(h(y_0)=0\) for some constant \(y_0\), then \(y(x)\equiv y_0\) is a \textbf{constant solution} of the original DE (since \(dy/dx=0\) and the right‑hand side becomes \(g(x)\cdot 0 = 0\)).  These constant solutions are \emph{not} captured by the separated integration and must be checked separately.
\end{attention}

\courssub{Key Points to Remember}

\begin{aretenir}[Separable Equations — Essentials]
\begin{itemize}
\item A first‑order DE is separable iff it can be written as \(\frac{dy}{dx}=g(x)h(y)\).
\item The separated form is \(\frac{1}{h(y)}\,dy = g(x)\,dx\) (or \(f(x)\,dx = g(y)\,dy\)).
\item Always check for constant solutions \(y=y_0\) where \(h(y_0)=0\) before dividing by \(h(y)\).
\item Common algebraic tricks: factorisation, \(e^{x+y}=e^xe^y\), \(\ln(xy)=\ln x+\ln y\), trigonometric identities.
\item Sums like \(x+y\) or mixed terms like \(x^2+y^2\) in denominators generally prevent separability.
\end{itemize}
\end{aretenir}

\courssub{Worked Example: Recognising and Solving a Separable DE}

\begin{exoresolu}[Recognise and solve \(\dfrac{dy}{dx} = \dfrac{2x}{y(1+x^2)}\)]
\textbf{Statement.}  Determine whether the differential equation
\[
\frac{dy}{dx} = \frac{2x}{y(1+x^2)}
\]
is separable.  If it is, find its general solution and identify any constant solutions.

\textbf{Detailed solution.}

\begin{enumerate}
\item \textbf{Write in explicit form.}  The equation is already given as \(dy/dx = F(x,y)\) with
\[
F(x,y) = \frac{2x}{y(1+x^2)}.
\]

\item \textbf{Factor \(F(x,y)\) as \(g(x)h(y)\).}  Observe that
\[
\frac{2x}{y(1+x^2)} = \underbrace{\frac{2x}{1+x^2}}_{g(x)}\;\underbrace{\frac{1}{y}}_{h(y)}.
\]
Thus the DE is separable with \(g(x)=\dfrac{2x}{1+x^2}\) and \(h(y)=\dfrac{1}{y}\).

\item \textbf{Check for constant solutions.}  \(h(y)=1/y = 0\) has no real solution, so no constant solutions are lost by division.

\item \textbf{Separate the variables.}
\[
\frac{1}{h(y)}\,dy = g(x)\,dx \quad\Longrightarrow\quad y\,dy = \frac{2x}{1+x^2}\,dx.
\]

\item \textbf{Integrate both sides.}
\[
\int y\,dy = \int \frac{2x}{1+x^2}\,dx.
\]
Left side: \(\frac{y^2}{2}\).  Right side: substitute \(u=1+x^2\), \(du=2x\,dx\) giving \(\int \frac{du}{u} = \ln|u| = \ln(1+x^2)\) (since \(1+x^2>0\)).  Hence
\[
\frac{y^2}{2} = \ln(1+x^2) + C,
\]
where \(C\) is an arbitrary constant.

\item \textbf{Express the general solution.}
\[
y^2 = 2\ln(1+x^2) + 2C \quad\Longrightarrow\quad y = \pm\sqrt{2\ln(1+x^2) + C'},
\]
with \(C'=2C\) an arbitrary constant.  The \(\pm\) reflects the two branches obtained when taking the square root.  The solution is defined for \(2\ln(1+x^2)+C' \ge 0\).

\item \textbf{Final answer.}  The DE is separable.  Its general solution is
\[
\boxed{y = \pm\sqrt{2\ln(1+x^2) + C}},\qquad C\in\mathbb{R},
\]
and there are no constant solutions.
\end{enumerate}
\end{exoresolu}

\courssub{Practice Set: Decide Separability}

\begin{exercice}[Separable or not?]
For each of the following differential equations, decide whether it is separable.  If it is, write it in the separated form \(f(x)\,dx = g(y)\,dy\).  If it is not, give a brief reason.

\begin{enumerate}
\item \(\dfrac{dy}{dx} = x^2 y^3\)
\item \(\dfrac{dy}{dx} = \sin(x+y)\)
\item \(\dfrac{dy}{dx} = \dfrac{e^x}{y^2+1}\)
\item \(\dfrac{dy}{dx} = \dfrac{x}{y} + \dfrac{y}{x}\)
\item \(\dfrac{dy}{dx} = \dfrac{xy}{x^2+y^2}\)
\item \(\dfrac{dy}{dx} = \dfrac{\cos x}{\sin y}\)
\item \(\dfrac{dy}{dx} = \dfrac{y^2-4}{x^2-9}\)
\item \(\dfrac{dy}{dx} = \dfrac{x+y}{x-y}\)
\item \(\dfrac{dy}{dx} = \dfrac{e^{x-y}}{x}\)
\item \(\dfrac{dy}{dx} = \dfrac{1}{x} + \dfrac{1}{y}\)
\end{enumerate}
\end{exercice}

\courssub{Answers and Brief Comments}

\begin{corrige}[Practice Set]
\begin{enumerate}
\item \textbf{Separable.}  \(g(x)=x^2,\; h(y)=y^3\).  Separated: \(\dfrac{1}{y^3}\,dy = x^2\,dx\).
\item \textbf{Not separable.}  \(\sin(x+y)\) cannot be factorised into a product of a function of \(x\) and a function of \(y\).
\item \textbf{Separable.}  \(g(x)=e^x,\; h(y)=\dfrac{1}{y^2+1}\).  Separated: \((y^2+1)\,dy = e^x\,dx\).
\item \textbf{Not separable.}  The sum \(\dfrac{x}{y}+\dfrac{y}{x} = \dfrac{x^2+y^2}{xy}\) mixes variables; no factorisation into \(g(x)h(y)\).
\item \textbf{Not separable.}  Denominator \(x^2+y^2\) prevents separation.
\item \textbf{Separable.}  \(g(x)=\cos x,\; h(y)=\dfrac{1}{\sin y}\).  Separated: \(\sin y\,dy = \cos x\,dx\).
\item \textbf{Separable.}  \(g(x)=\dfrac{1}{x^2-9},\; h(y)=y^2-4\).  Separated: \(\dfrac{1}{y^2-4}\,dy = \dfrac{1}{x^2-9}\,dx\).
\item \textbf{Not separable.}  The quotient \(\dfrac{x+y}{x-y}\) is not a product of separate functions.
\item \textbf{Separable.}  \(e^{x-y}=e^x e^{-y}\).  \(g(x)=\dfrac{e^x}{x},\; h(y)=e^{-y}\).  Separated: \(e^{y}\,dy = \dfrac{e^x}{x}\,dx\).
\item \textbf{Not separable.}  Sum of a function of \(x\) and a function of \(y\) cannot be written as a product.
\end{enumerate}
\end{corrige}

\courssub{Looking Ahead}

Now that we can confidently recognise separable equations, the next section will show how to carry out the integration step to obtain the \cle{general solution}, and how to use initial conditions to pick out a \cle{particular solution}.  We will also explore a variety of real‑world applications — from population growth in the Adamawa region to cooling of a cup of Cameroonian coffee — where separable differential equations provide exact mathematical models.

\coursec{Separation of Variables}

This section covers \textbf{Lessons 63 and 64} of the official syllabus: the technique of separating variables and integrating to obtain the general solution of a first order separable differential equation. In the previous section we learned to \emph{recognise} a separable equation: one that can be written in the form $\dfrac{dy}{dx} = g(x)\,h(y)$, where the right-hand side splits into a factor depending only on $x$ and a factor depending only on $y$. Recognition, however, is only the first step. We now learn the central skill of this chapter: how to \emph{solve} such an equation by \cle{separating the variables} — gathering everything involving $y$ with $dy$ on one side, and everything involving $x$ with $dx$ on the other — and then \emph{integrating both sides}. This single idea, simple as it looks, unlocks an enormous family of differential equations, from population growth to radioactive decay and Newton's law of cooling.

\courssub{The Formal Separation Procedure}

\begin{definition}[Separation of variables]
Let
\[ \frac{dy}{dx} = g(x)\,h(y) \]
be a first order separable differential equation, where $g$ is continuous on an interval $I$ and $h$ is continuous on an interval $J$. On any interval where $h(y) \neq 0$, we may divide both sides by $h(y)$ and write the equation in \textbf{separated form}:
\[ \frac{1}{h(y)}\,dy = g(x)\,dx. \]
Integrating both sides with respect to their own variable gives
\[ \int \frac{dy}{h(y)} = \int g(x)\,dx + C, \]
that is, an implicit relation $F(y) = G(x) + C$ between $x$ and $y$, where $F'(y) = 1/h(y)$ and $G'(x) = g(x)$.
\end{definition}

\textbf{Why is this legitimate?}\par
A careful justification uses the chain rule. Suppose $F$ is an antiderivative of $\dfrac{1}{h}$, so $F'(y) = \dfrac{1}{h(y)}$. If $y$ is a function of $x$ satisfying $\dfrac{dy}{dx} = g(x)h(y)$, then
\[ \frac{d}{dx}\big[F(y)\big] = F'(y)\,\frac{dy}{dx} = \frac{1}{h(y)}\cdot g(x)h(y) = g(x). \]
Hence $F(y)$ is an antiderivative of $g(x)$, so $F(y) = G(x) + C$ where $G'(x) = g(x)$. This is exactly what the symbolic manipulation ``$\dfrac{dy}{h(y)} = g(x)\,dx$, then integrate'' produces. The Leibniz notation $\dfrac{dy}{dx}$ is doing real work for us here: treating $dy$ and $dx$ as separable quantities is a reliable shortcut, and the chain rule guarantees that the shortcut is sound.

\begin{methode}[Separation of variables — the standard routine]
To solve $\dfrac{dy}{dx} = g(x)h(y)$:
\begin{enumerate}
\item \textbf{Write the equation in separable form.} Identify $g(x)$ and $h(y)$ explicitly.
\item \textbf{Isolate the differentials.} Divide by $h(y)$ (noting the values of $y$ for which $h(y)=0$) and multiply by $dx$ to obtain $\dfrac{dy}{h(y)} = g(x)\,dx$.
\item \textbf{Integrate both sides.} Compute $\displaystyle\int \frac{dy}{h(y)}$ and $\displaystyle\int g(x)\,dx$.
\item \textbf{Add one constant of integration.} A single constant $C$ on one side is enough, since two constants $C_1, C_2$ merge into $C = C_2 - C_1$.
\item \textbf{Make $y$ the subject if possible}, and state the constant solutions (if any) found in step 1.
\end{enumerate}
\end{methode}

\begin{popfigure}
\begin{center}
\begin{tikzpicture}[
box/.style={draw=popblue, fill=popblueL, rounded corners=2pt, align=center, text width=4.1cm, minimum height=1.15cm, font=\small},
arr/.style={-{Stealth[length=3mm]}, very thick, poporange},
lbl/.style={font=\footnotesize\itshape, text=popdark}]
\node[box] (start) at (0,0) {$\dfrac{dy}{dx}=\dfrac{xy}{1+x^{2}}$};
\node[box] (sep) at (6.9,0) {$\dfrac{dy}{y}=\dfrac{x}{1+x^{2}}\,dx$};
\node[box] (int) at (13.8,0) {$\displaystyle\int\frac{dy}{y}=\int\frac{x}{1+x^{2}}\,dx$};
\node[box] (res) at (13.8,-4.5) {$\ln|y|=\tfrac12\ln(1+x^{2})+C$};
\node[box] (exp) at (6.9,-4.5) {$|y|=A\sqrt{1+x^{2}}$};
\node[box] (fin) at (0,-4.5) {$y=K\sqrt{1+x^{2}},\ K\in\mathbb{R}$};
\draw[arr] (start) -- (sep) node[lbl, midway, above]{divide by $y$, multiply by $dx$};
\draw[arr] (sep) -- (int) node[lbl, midway, above]{integrate both sides};
\draw[arr] (int) -- (res) node[lbl, midway, right]{evaluate integrals};
\draw[arr] (res) -- (exp) node[lbl, midway, above]{exponentiate, $A=e^{C}>0$};
\draw[arr] (exp) -- (fin) node[lbl, midway, above]{absorb sign into $K$};
\end{tikzpicture}
\end{center}
\legende{Step-by-step separation of variables for $\dfrac{dy}{dx}=\dfrac{xy}{1+x^{2}}$.}
\end{popfigure}

\courssub{Do Not Forget the Constant (Equilibrium) Solutions}

When we divided by $h(y)$ in the very first step, we silently assumed $h(y) \neq 0$. But what happens at the values of $y$ where $h(y) = 0$? If $h(y_0) = 0$ for some constant $y_0$, then the constant function $y(x) = y_0$ has derivative $\dfrac{dy}{dx} = 0$, and the right-hand side is $g(x)h(y_0) = g(x)\cdot 0 = 0$. So the equation is satisfied identically.

\begin{propriete}[Constant solutions]
If $h(y_0) = 0$, then the constant function $y = y_0$ is a solution of $\dfrac{dy}{dx} = g(x)h(y)$ on the whole interval where $g$ is defined. These are called \cle{constant solutions} or \cle{equilibrium solutions}. (The verification is examinable: substitute and check both sides vanish.)
\end{propriete}

\begin{attention}[The incomplete general solution]
Dividing by $h(y)$ without checking where $h(y) = 0$ is the most common way to \emph{lose} solutions. For example, in $\dfrac{dy}{dx} = y^2\sin x$, dividing by $y^2$ throws away the solution $y = 0$, which is perfectly valid. Always ask: \emph{``Can $h(y)$ be zero? If so, is that constant a solution?''} Some of these constant solutions are recovered from the general formula for a special value of the constant (e.g. $K = 0$ in $y = K\sqrt{1+x^2}$); others, like $y = 0$ for $\frac{dy}{dx}=y^2\sin x$, are \textbf{singular solutions} that no choice of $C$ will ever produce.
\end{attention}

\courssub{Integration: From Separated Form to Implicit Relation}

After separation, each side is integrated with respect to its own variable. The result is a relation
\[ F(y) = G(x) + C, \]
which usually defines $y$ only \emph{implicitly} as a function of $x$. Two practical points deserve emphasis.

\begin{aretenir}[Handling $\int \frac{dy}{y}$]
Whenever the integral produces a logarithm, write the absolute value: $\displaystyle\int \frac{dy}{y} = \ln|y| + c$. Do \textbf{not} drop the modulus signs. The sign of $y$ is handled later, cleanly, by allowing the arbitrary constant to be positive or negative when you exponentiate: $|y| = e^{C}\cdot(\ldots)$ becomes $y = K\cdot(\ldots)$ with $K \in \mathbb{R}\setminus\{0\}$, and $K = 0$ is then checked separately.
\end{aretenir}

\begin{exoresolu}[Worked example 1: an implicit solution]
Solve the differential equation $\dfrac{dy}{dx} = \dfrac{xy}{1+x^{2}}$.
\tcblower
\textbf{Step 1 — Constant solutions.} Here $h(y) = y$, which vanishes at $y = 0$. The constant function $y = 0$ is indeed a solution (both sides equal $0$).

\textbf{Step 2 — Separate.} For $y \neq 0$:
\[ \frac{dy}{y} = \frac{x}{1+x^{2}}\,dx. \]

\textbf{Step 3 — Integrate.} The left side gives $\ln|y|$. For the right side, notice that the numerator $x$ is (up to a factor $2$) the derivative of $1+x^2$:
\[ \int \frac{x}{1+x^{2}}\,dx = \frac{1}{2}\ln(1+x^{2}) + c. \]
Hence
\[ \ln|y| = \frac{1}{2}\ln(1+x^{2}) + C. \]
This is the \textbf{implicit solution}. (Note: $1+x^2 > 0$ always, so no modulus is needed there.)

\textbf{Step 4 — Make $y$ explicit.} Exponentiating,
\[ |y| = e^{C}\,(1+x^{2})^{1/2}. \]
Writing $K = \pm e^{C}$ (any non-zero real), and observing that $K = 0$ restores the constant solution $y=0$ lost in Step 2, the general solution is
\[ \boxed{\,y = K\sqrt{1+x^{2}},\quad K \in \mathbb{R}.\,} \]

\textbf{Check.} If $y = K\sqrt{1+x^2}$, then $\dfrac{dy}{dx} = \dfrac{Kx}{\sqrt{1+x^2}} = \dfrac{x}{1+x^2}\cdot K\sqrt{1+x^2} = \dfrac{xy}{1+x^2}$. \checkmark
\end{exoresolu}

\courssub{From Implicit to Explicit Solutions}

The relation $F(y) = G(x) + C$ is a perfectly acceptable answer — it is the \cle{general solution in implicit form}. But whenever algebra permits, we solve it for $y$ to obtain the \cle{explicit solution} $y = \varphi(x)$. This is possible when $F$ can be inverted: typical tools are exponentiation (to undo $\ln$), taking reciprocals, or taking roots.

\begin{propriete}[Branches of an implicit solution]
An implicit relation $F(y) = G(x) + C$ may define \textbf{several branches} $y = \varphi(x)$ (for instance, $y^2 = x + C$ gives $y = \sqrt{x+C}$ and $y = -\sqrt{x+C}$). Each branch must be checked against the \emph{original} differential equation: only branches that satisfy it, on an interval where they are differentiable, are genuine solutions. (This verification is examinable.)
\end{propriete}

\begin{attention}[Extraneous solutions from algebraic manipulation]
Operations performed \emph{after} integration can smuggle in false solutions. Squaring both sides turns $A = B$ into $A^2 = B^2$, which also admits $A = -B$; taking square roots introduces a $\pm$ that must be analysed, not ignored. \textbf{Rule of thumb:} after any non-reversible manipulation, substitute your final formula back into the original equation and check the sign.
\end{attention}

\begin{exoresolu}[Worked example 2: an explicit solution]
Find the general solution of $\dfrac{dy}{dx} = y^{2}\sin x$.
\tcblower
\textbf{Step 1 — Constant solutions.} $h(y) = y^2 = 0$ when $y = 0$. The function $y = 0$ satisfies the equation: it is a (singular) solution.

\textbf{Step 2 — Separate.} For $y \neq 0$:
\[ \frac{dy}{y^{2}} = \sin x\, dx. \]

\textbf{Step 3 — Integrate.}
\[ \int y^{-2}\,dy = \int \sin x\, dx \quad\Longrightarrow\quad -\frac{1}{y} = -\cos x + C. \]

\textbf{Step 4 — Make $y$ explicit.} Multiply by $-1$: $\dfrac{1}{y} = \cos x - C$. Renaming the arbitrary constant as $c = -C$ (still an arbitrary real constant),
\[ \frac{1}{y} = \cos x + c \quad\Longrightarrow\quad \boxed{\,y = \frac{1}{\cos x + c}\,}. \]

\textbf{Full answer.} The general solution is $y = \dfrac{1}{\cos x + c}$ ($c \in \mathbb{R}$), together with the singular solution $y = 0$, which is not obtained for any value of $c$.

\textbf{Check.} If $y = (\cos x + c)^{-1}$, then $\dfrac{dy}{dx} = -(\cos x + c)^{-2}(-\sin x) = y^{2}\sin x$. \checkmark
\end{exoresolu}

\courssub{Domain Considerations}

A solution of a differential equation is, strictly speaking, a \emph{function defined on an interval}. The formula $y = \dfrac{1}{\cos x + c}$ from Worked Example 2 is meaningless where $\cos x + c = 0$: the solution curve blows up there. Hence:

\begin{aretenir}[Maximal interval of a solution]
\begin{itemize}
\item The solution must be \textbf{continuous and differentiable} on its interval of definition; it cannot jump across a vertical asymptote.
\item The interval must lie inside the domain where the original equation makes sense (where $g$ is continuous and $h(y)\neq 0$ along the solution).
\item The \cle{maximal solution} is the solution defined on the \textbf{largest possible interval} containing the initial point. For $y = \dfrac{1}{\cos x + c}$ with $c = 2$, the denominator never vanishes (since $\cos x \geq -1 > -2$), so the maximal interval is the whole of $\mathbb{R}$; with $c = 0$, the denominator vanishes at $x = \frac{\pi}{2} + k\pi$, and each maximal solution lives on a single interval $\left(-\frac{\pi}{2}+k\pi,\ \frac{\pi}{2}+k\pi\right)$.
\end{itemize}
\end{aretenir}

This is why, at GCE Advanced Level, an answer such as ``$y = \dfrac{1}{\cos x + c}$'' should ideally be accompanied by a remark on where it is valid — especially once an initial condition (next section) picks out one particular branch and one particular interval.

\begin{savaistu}[Why ``separation'' is everywhere]
Separable equations are the workhorse of applied mathematics. Radioactive decay $\dfrac{dN}{dt} = -\lambda N$, Newton's law of cooling $\dfrac{dT}{dt} = -k(T - T_{\text{env}})$, and the growth of a savings account at continuous interest $\dfrac{dM}{dt} = rM$ are all separable — and all solved by exactly the routine of this section. Master these four steps now, and Section 5's applications will feel like revision.
\end{savaistu}

\begin{exercice}[Practice]
Solve the following equations by separation of variables, giving the general solution and stating any constant solutions:
\begin{enumerate}
\item $\dfrac{dy}{dx} = \dfrac{x}{y}$ \quad (leave the answer in implicit form);
\item $\dfrac{dy}{dx} = y\cos x$, and state the maximal interval of the solution through $(0, 2)$;
\item $\dfrac{dy}{dx} = \dfrac{y^{2}}{x}$ for $x > 0$.
\end{enumerate}
\end{exercice}

\begin{corrige}[Exercises 1–3]
\textbf{1.} Separate: $y\,dy = x\,dx$. Integrate: $\frac{y^{2}}{2} = \frac{x^{2}}{2} + C$, i.e. $y^{2} - x^{2} = c$ (hyperbolas). No constant solution since $h(y) = \frac{1}{y}$ never vanishes — but note $y = 0$ is excluded as the original equation is undefined there.

\textbf{2.} $y = 0$ is a constant solution. For $y \neq 0$: $\frac{dy}{y} = \cos x\,dx$, so $\ln|y| = \sin x + C$, giving $y = Ke^{\sin x}$, $K \in \mathbb{R}$ (including $K=0$). Through $(0,2)$: $K = 2$, so $y = 2e^{\sin x}$, defined and smooth on the whole of $\mathbb{R}$ — the maximal interval is $\mathbb{R}$.

\textbf{3.} $y = 0$ is a singular solution. For $y \neq 0$: $\frac{dy}{y^{2}} = \frac{dx}{x}$, so $-\frac{1}{y} = \ln x + C$, i.e. $y = -\dfrac{1}{\ln x + C}$, valid on intervals where $\ln x + C \neq 0$ within $x > 0$.
\end{corrige}

\coursec{Integration to Obtain General Solutions}

In the previous section we learned how to \emph{separate} the variables of a first order differential equation: we transformed an equation of the form $\dfrac{dy}{dx}=g(x)\,h(y)$ into the form
\[
\frac{1}{h(y)}\,dy = g(x)\,dx,
\]
where all the $y$'s live on the left and all the $x$'s live on the right. Separation, however, is only half of the journey. The equation still contains differentials; we have not yet found $y$ as a function of $x$. The decisive step is now to \textbf{integrate both sides}. This section explains exactly how to do that, how to handle the constant of integration, how to simplify the result (especially when logarithms appear), and how to check that the answer is correct.

\courssub{The Fundamental Step: Integrating Both Sides}

\begin{definition}[General solution by direct integration]
Once the variables of $\dfrac{dy}{dx}=g(x)h(y)$ (with $h(y)\neq 0$) have been separated, the \cle{general solution} is obtained by integrating each side with respect to its own variable:
\[
\int \frac{1}{h(y)}\,dy \;=\; \int g(x)\,dx \;+\; C,
\]
where $C$ is an arbitrary constant of integration.
\end{definition}

Why is this legitimate? If two functions of $x$ have the same derivative, they differ by a constant. Writing $H(y)=\displaystyle\int \frac{dy}{h(y)}$, the separated equation says $\dfrac{d}{dx}\big[H(y)\big]=g(x)$: indeed, by the chain rule,
\[
\frac{d}{dx}\big[H(y)\big]=\frac{dH}{dy}\cdot\frac{dy}{dx}=\frac{1}{h(y)}\cdot g(x)h(y)=g(x).
\]
Integrating both sides with respect to $x$ gives $H(y)=\displaystyle\int g(x)\,dx + C$, which is exactly the formula in the definition.

\begin{methode}[Integrating after separation]
\begin{enumerate}
\item Separate the variables completely: $\dfrac{1}{h(y)}\,dy=g(x)\,dx$.
\item Integrate the left side \textbf{with respect to $y$} and the right side \textbf{with respect to $x$}.
\item Add \textbf{one single} arbitrary constant $C$ (it is enough to put it on one side, since two constants $C_1,C_2$ combine into $C=C_2-C_1$).
\item Simplify: combine logarithms, exponentiate if useful, and absorb constants.
\item State the general solution, in implicit or explicit form as required, indicating the values of the variable for which it is valid.
\end{enumerate}
\end{methode}

\begin{attention}[Do not forget the constant]
Omitting the constant of integration does \textbf{not} give the general solution — it gives only \emph{one} particular solution. Since a first order equation has a one-parameter family of solutions, the general solution must contain exactly one arbitrary constant. In the GCE examination, a missing $+C$ costs marks every time.
\end{attention}

\courssub{Integration Techniques You Will Need}

Separable equations reduce the problem to two ordinary integrals, so all the integration skills of the Lower Sixth and Upper Sixth course are mobilised. Here is a rapid review of the tools that appear most often.

\textbf{Standard integrals.}
\begin{center}
\begin{tabularx}{\linewidth}{|l|X|}
\hline
\rowcolor{popblueL} \textbf{Integral} & \textbf{Result} \\
\hline
$\displaystyle\int x^n\,dx\ (n\neq -1)$ & $\dfrac{x^{n+1}}{n+1}+C$ \\
\hline
$\displaystyle\int \frac{1}{x}\,dx$ & $\ln|x|+C$ \\
\hline
$\displaystyle\int e^{ax}\,dx$ & $\dfrac{1}{a}e^{ax}+C$ \\
\hline
$\displaystyle\int \frac{dx}{a^2+x^2}$ & $\dfrac{1}{a}\arctan\dfrac{x}{a}+C$ \\
\hline
$\displaystyle\int \frac{dx}{\sqrt{a^2-x^2}}$ & $\arcsin\dfrac{x}{a}+C$ \\
\hline
$\displaystyle\int \frac{f'(x)}{f(x)}\,dx$ & $\ln|f(x)|+C$ \\
\hline
\end{tabularx}
\end{center}

\textbf{Substitution.} When the integrand has the form $f'(x)\cdot F(f(x))$, substitute $u=f(x)$. For example, $\displaystyle\int \frac{2x}{1+x^2}\,dx$: with $u=1+x^2$, $du=2x\,dx$, the integral becomes $\displaystyle\int \frac{du}{u}=\ln(1+x^2)+C$ (no absolute value is needed since $1+x^2>0$).

\textbf{Partial fractions.} Rational expressions such as $\dfrac{1}{y(1-y)}$ (very common in population models) must be decomposed first:
\[
\frac{1}{y(1-y)}=\frac{1}{y}+\frac{1}{1-y},
\qquad\text{so}\qquad
\int\frac{dy}{y(1-y)}=\ln|y|-\ln|1-y|+C=\ln\left|\frac{y}{1-y}\right|+C.
\]

\begin{exemplebox}[Warm-up integrals]
\begin{align*}
\int \frac{y\,dy}{1+y^2} &= \tfrac12\ln(1+y^2)+C &&\text{(substitution } u=1+y^2\text{)}\\
\int \frac{dx}{x^2+4} &= \tfrac12\arctan\frac{x}{2}+C &&\text{(standard form with } a=2\text{)}
\end{align*}
\end{exemplebox}

\courssub{A Complete Model Example: $\dfrac{dy}{dx}=\dfrac{2x}{y}$}

This example is the archetype of the method; master it and the whole section follows.

\begin{exoresolu}[General solution of $\dfrac{dy}{dx}=\dfrac{2x}{y}$]
Find the general solution of the differential equation $\dfrac{dy}{dx}=\dfrac{2x}{y}$.
\tcblower
\textbf{Step 1 — Separate.} Multiply both sides by $y\,dx$:
\[
y\,dy = 2x\,dx.
\]
\textbf{Step 2 — Integrate both sides.}
\[
\int y\,dy = \int 2x\,dx \quad\Longrightarrow\quad \frac{y^2}{2}=x^2+C_1.
\]
\textbf{Step 3 — Simplify.} Multiply by $2$ and rename the constant $C=2C_1$:
\[
\[\boxed{\,y^2 = 2x^2 + C\,}\]
\]
This is the \textbf{general solution in implicit form}. Solving for $y$ gives the explicit form $y=\pm\sqrt{2x^2+C}$ (valid where $2x^2+C\geq 0$): each choice of sign and of $C$ gives one particular solution.
\end{exoresolu}

\begin{popfigure}
\begin{tikzpicture}[
  solstep/.style={draw=popblue, thick, rounded corners=3pt, fill=popblueL, text width=5.6cm, align=center, inner sep=5pt, font=\small},
  arr/.style={-{Stealth[length=2.5mm]}, thick, poporange},
  lab/.style={font=\footnotesize\itshape, text=popdark}]
\node[solstep] (s1) at (0,0) {$\dfrac{dy}{dx}=\dfrac{2x}{y}$};
\node[solstep] (s2) at (0,-1.7) {$y\,dy = 2x\,dx$};
\node[solstep] (s3) at (0,-3.4) {$\displaystyle\int y\,dy=\int 2x\,dx$};
\node[solstep] (s4) at (0,-5.1) {$\dfrac{y^2}{2}=x^2+C_1$};
\node[solstep] (s5) at (0,-6.8) {$y^2=2x^2+C$};
\draw[arr] (s1) -- (s2) node[midway, right, lab]{separate};
\draw[arr] (s2) -- (s3) node[midway, right, lab]{integrate both sides};
\draw[arr] (s3) -- (s4) node[midway, right, lab]{evaluate, add one constant};
\draw[arr] (s4) -- (s5) node[midway, right, lab]{simplify: $C=2C_1$};
\end{tikzpicture}
\legende{Integration steps for $\frac{dy}{dx}=\frac{2x}{y}$, from the equation to the general solution.}
\end{popfigure}

The family of solutions can also be visualised: different values of $C$ give different curves, all sharing the same ``slope field'' defined by the differential equation.

\begin{popfigure}
\begin{center}
\begin{tikzpicture}
\begin{axis}[
  width=9cm, height=6.5cm,
  axis lines=middle, xlabel={$x$}, ylabel={$y$},
  xmin=-2.5, xmax=2.5, ymin=-4, ymax=4,
  domain=-2.3:2.3, samples=120,
  legend style={font=\footnotesize, at={(0.02,0.98)}, anchor=north west}]
\addplot[popblue, thick] {sqrt(max(0,2*x^2+1))}; \addlegendentry{$C=1$}
\addplot[poporange, thick] {sqrt(max(0,2*x^2+4))}; \addlegendentry{$C=4$}
\addplot[poppurple, thick] {-sqrt(max(0,2*x^2+1))};
\addplot[popgreen, thick] {-sqrt(max(0,2*x^2+4))};
\end{axis}
\end{tikzpicture}
\end{center}
\legende{Some particular solutions $y=\pm\sqrt{2x^2+C}$ of $\frac{dy}{dx}=\frac{2x}{y}$. The general solution is the whole family, parametrised by $C$.}
\end{popfigure}

\courssub{Working with Logarithms: Absolute Values and the $\pm$ Sign}

Many separable equations produce logarithms, because $\displaystyle\int \frac{dy}{y}=\ln|y|$. Two technical points then arise constantly at the GCE.

\textbf{Combining logarithms before exponentiating.} Always use the laws $\ln A+\ln B=\ln(AB)$ and $n\ln A=\ln(A^n)$ to gather all logarithms on one side \emph{before} removing them.

\begin{methode}[Removing logarithms]
\begin{enumerate}
\item Gather all logarithmic terms on one side, constants on the other.
\item Combine them into a \emph{single} logarithm.
\item Exponentiate both sides: $\ln|y|=F(x)+C_1 \Rightarrow |y|=e^{F(x)+C_1}=e^{C_1}e^{F(x)}$.
\item Remove the absolute value: $y=\pm e^{C_1}e^{F(x)}$, and write $A=\pm e^{C_1}$, an arbitrary non-zero constant.
\end{enumerate}
\end{methode}

\begin{attention}[The $\pm$ after exponentiation]
From $\ln|y|=F(x)+C_1$ we get $|y|=e^{C_1}e^{F(x)}$, hence $y=\pm e^{C_1}e^{F(x)}$. The factor $\pm e^{C_1}$ is just \emph{some} non-zero constant: we rename it $A$ (or $K$). Writing $y=e^{C_1}e^{F(x)}$ and forgetting the $\pm$ silently throws away half of the solutions (all the negative ones). Note also that $y=0$, excluded when we divided by $y$, is often recovered by allowing $A=0$.
\end{attention}

\begin{exoresolu}[A logarithmic case: $\dfrac{dy}{dx}=\dfrac{y}{x}$]
Find the general solution of $\dfrac{dy}{dx}=\dfrac{y}{x}$ for $x>0$, $y\neq 0$.
\tcblower
\textbf{Step 1 — Separate.}
\[
\frac{dy}{y}=\frac{dx}{x}.
\]
\textbf{Step 2 — Integrate.}
\[
\ln|y|=\ln|x|+C_1.
\]
\textbf{Step 3 — Combine and exponentiate.} Since $x>0$, $|x|=x$:
\[
\ln|y|-\ln x = C_1 \;\Longrightarrow\; \ln\left|\frac{y}{x}\right|=C_1 \;\Longrightarrow\; \left|\frac{y}{x}\right|=e^{C_1}.
\]
\textbf{Step 4 — Absorb the $\pm$.} $\dfrac{y}{x}=\pm e^{C_1}=A$, so
\[
\[\boxed{\,y=Ax\quad (A\in\mathbb{R},\,A\neq 0)\,}\]
\]
and allowing $A=0$ also recovers the solution $y=0$ lost during separation. The general solution is the family of all straight lines through the origin.
\end{exoresolu}

\begin{propriete}[Implicit and explicit forms are equivalent]
The general solution may be left in \cle{implicit form} (e.g. $y^2=2x^2+C$) or \cle{explicit form} ($y=\pm\sqrt{2x^2+C}$): both describe exactly the same family of curves. Choose the form requested by the question; if nothing is specified, the simplest implicit form is perfectly acceptable. (Not examinable as a proof, but the equivalence must be understood.)
\end{propriete}

\courssub{One Constant Only: Absorbing and Renaming}

Students are often puzzled that constants seem to ``change'' during simplification. The principle is simple: \emph{any fixed function of an arbitrary constant is again an arbitrary constant}, so we rename it to keep the answer clean.

\begin{exemplebox}[Renaming constants]
From $\dfrac{y^2}{2}=x^2+C_1$ we wrote $y^2=2x^2+C$ with $C=2C_1$. From $\ln|y|=3\ln|x|+C_1$ we get $|y|=e^{C_1}|x|^3$, i.e. $y=Ax^3$ with $A=\pm e^{C_1}$. In each case the final answer contains exactly \textbf{one} arbitrary constant — the signature of a first order equation.
\end{exemplebox}

\begin{attention}[Trap: two constants]
Never write $\displaystyle\int y\,dy=\int 2x\,dx$ as $\frac{y^2}{2}+C_1=x^2+C_2$ and leave both: they merge into $C=C_2-C_1$. Two independent constants would wrongly suggest a second order equation.
\end{attention}

\courssub{A Harder Case: Partial Fractions in Action}

Some GCE questions require a partial fraction decomposition before the integration can be carried out. The following worked exercise shows the complete chain.

\begin{exoresolu}[Using partial fractions: $\dfrac{dy}{dx}=\dfrac{y(2-y)}{x}$ for $x>0$, $0<y<2$]
Find the general solution of $\dfrac{dy}{dx}=\dfrac{y(2-y)}{x}$, giving $y$ explicitly in terms of $x$.
\tcblower
\textbf{Step 1 — Separate.}
\[
\frac{dy}{y(2-y)}=\frac{dx}{x}.
\]
\textbf{Step 2 — Decompose.} Write $\dfrac{1}{y(2-y)}=\dfrac{A}{y}+\dfrac{B}{2-y}$. Multiplying by $y(2-y)$: $1=A(2-y)+By$. Setting $y=0$ gives $A=\tfrac12$; setting $y=2$ gives $B=\tfrac12$. Hence
\[
\frac{1}{y(2-y)}=\frac{1}{2y}+\frac{1}{2(2-y)}.
\]
\textbf{Step 3 — Integrate.} Since $0<y<2$ and $x>0$, all quantities inside the logarithms are positive:
\[
\frac12\ln y-\frac12\ln(2-y)=\ln x + C_1
\;\Longrightarrow\;
\ln\frac{y}{2-y}=2\ln x + 2C_1=\ln(x^2)+2C_1.
\]
\textbf{Step 4 — Exponentiate and solve for $y$.}
\[
\frac{y}{2-y}=e^{2C_1}x^2=Kx^2 \quad (K>0)
\;\Longrightarrow\;
y=Kx^2(2-y)
\;\Longrightarrow\;
y(1+Kx^2)=2Kx^2.
\]
\[
\[\boxed{\,y=\frac{2Kx^2}{1+Kx^2}\,}\]
\]
Note that $0<y<2$ for every $K>0$, consistent with the hypothesis. The constant solutions $y=0$ and $y=2$, excluded when dividing by $y(2-y)$, are singular solutions of the equation.
\end{exoresolu}

\courssub{Verification: Differentiate to Recover the Equation}

A powerful habit, which costs one minute and can save the whole question: \textbf{differentiate your answer implicitly} and check that you recover the original differential equation.

\begin{exemplebox}[Checking $y^2=2x^2+C$]
Differentiate with respect to $x$:
\[
2y\frac{dy}{dx}=4x \quad\Longrightarrow\quad \frac{dy}{dx}=\frac{2x}{y},
\]
which is exactly the given equation. The constant $C$ has disappeared, as it must, confirming that \emph{every} member of the family satisfies the equation. \checkmark
\end{exemplebox}

\begin{exoresolu}[Full worked exercise: $\dfrac{dy}{dx}=\dfrac{x}{y(1+x^2)}$]
Find the general solution of $\dfrac{dy}{dx}=\dfrac{x}{y(1+x^2)}$, expressing $y$ explicitly in terms of $x$, and verify your answer.
\tcblower
\textbf{Step 1 — Separate.} $y\,dy=\dfrac{x}{1+x^2}\,dx$.

\textbf{Step 2 — Integrate.} On the right, substitute $u=1+x^2$, $du=2x\,dx$:
\[
\int y\,dy=\int\frac{x}{1+x^2}\,dx
\;\Longrightarrow\;
\frac{y^2}{2}=\frac12\ln(1+x^2)+C_1.
\]
(Note $1+x^2>0$, so no absolute value is needed.)

\textbf{Step 3 — Simplify.} Multiply by $2$ and set $C=2C_1$:
\[
y^2=\ln(1+x^2)+C,
\qquad\text{hence}\qquad
\boxed{\,y=\pm\sqrt{\ln(1+x^2)+C}\,}.
\]

\textbf{Step 4 — Verify.} Differentiate $y^2=\ln(1+x^2)+C$:
\[
2y\frac{dy}{dx}=\frac{2x}{1+x^2}
\;\Longrightarrow\;
\frac{dy}{dx}=\frac{x}{y(1+x^2)}. \quad\checkmark
\]
\end{exoresolu}

\begin{exercice}[]
Find the general solution of each equation, then verify by implicit differentiation:
\begin{enumerate}
\item $\dfrac{dy}{dx}=3x^2 y$;
\item $\dfrac{dy}{dx}=\dfrac{1}{y(4+x^2)}$.
\end{enumerate}
\end{exercice}

\begin{corrige}[Exercise 1]
\textbf{1.} Separate: $\dfrac{dy}{y}=3x^2\,dx$. Integrate: $\ln|y|=x^3+C_1$. Exponentiate and absorb the $\pm$: $y=Ae^{x^3}$, $A\in\mathbb{R}$. Check: $\dfrac{dy}{dx}=3x^2\cdot Ae^{x^3}=3x^2y$. \checkmark

\textbf{2.} Separate: $y\,dy=\dfrac{dx}{4+x^2}$. Integrate using the standard arctangent form with $a=2$: $\dfrac{y^2}{2}=\dfrac12\arctan\dfrac{x}{2}+C_1$, so $y^2=\arctan\dfrac{x}{2}+C$. Check: differentiating, $2y\dfrac{dy}{dx}=\dfrac{1}{1+(x/2)^2}\cdot\dfrac12=\dfrac{2}{4+x^2}$, giving $\dfrac{dy}{dx}=\dfrac{1}{y(4+x^2)}$. \checkmark
\end{corrige}

\begin{aretenir}[Key points of this section]
\begin{itemize}
\item After separation, integrate each side with respect to its own variable: $\displaystyle\int\frac{dy}{h(y)}=\int g(x)\,dx+C$.
\item Add \textbf{one} arbitrary constant only; rename it freely when simplifying ($C=2C_1$, $A=\pm e^{C_1}$, \dots).
\item Use the full integration toolbox: standard forms, substitution, partial fractions.
\item Combine logarithms \emph{before} exponentiating; from $\ln|y|=F(x)+C_1$ deduce $y=Ae^{F(x)}$, the $\pm$ being absorbed into $A$.
\item Implicit and explicit forms are equivalent; give the form requested, and state the domain of validity when relevant.
\item Always verify: differentiate the solution implicitly and recover the original equation.
\end{itemize}
\end{aretenir}

\coursec{Particular Solutions and Applications}

In the previous section we learned how to separate the variables of a first order differential equation and integrate both sides to obtain the \emph{general solution}, which always contains an arbitrary constant $C$. A general solution is really a \emph{family} of curves: each value of $C$ gives one curve. In real problems, however, we are interested in \emph{one} specific situation: the population of \emph{this} village starting from a known size, the temperature of \emph{this} pot of ndolé measured at a known instant. The extra piece of information that selects one curve from the family is called an \emph{initial condition}, and the curve it selects is the \emph{particular solution}. This final section shows how to find particular solutions and how separable equations model growth, decay, cooling, mixing and electric circuits.

\courssub{From the General Solution to the Particular Solution}

\begin{definition}[Initial condition and particular solution]
An \cle{initial condition} for the equation $\dfrac{dy}{dx}=f(x)g(y)$ is a requirement of the form $y(x_0)=y_0$. A \cle{particular solution} is the member of the family of general solutions whose constant $C$ is chosen so that the initial condition is satisfied. The pair
\[
\frac{dy}{dx}=f(x)g(y), \qquad y(x_0)=y_0
\]
is called an \cle{initial value problem} (IVP).
\end{definition}

\begin{methode}[Finding a particular solution]
\begin{enumerate}
\item \textbf{Write the general solution} by separating variables and integrating (do not forget $+C$).
\item \textbf{Plug in the initial values}: replace $x$ by $x_0$ and $y$ by $y_0$ in the general solution.
\item \textbf{Solve for $C$} (this is usually a simple linear or logarithmic equation).
\item \textbf{Write the final solution} explicitly, solving for $y$ in terms of $x$ if possible, and state its domain.
\end{enumerate}
\end{methode}

\begin{propriete}[Existence and uniqueness (Picard--Lindelöf)]
If $f(x,y)$ and $\dfrac{\partial f}{\partial y}$ are continuous in a rectangle containing $(x_0,y_0)$, then the initial value problem $\dfrac{dy}{dx}=f(x,y)$, $y(x_0)=y_0$ has a \textbf{unique} solution on some interval around $x_0$. (Statement only; the proof is \emph{not} examinable at GCE level.)
\end{propriete}

This property is reassuring: when you find one solution satisfying the initial condition, you have found \emph{the} solution. There is no other curve to look for.

\begin{exoresolu}[A basic particular solution]
Solve the initial value problem $\dfrac{dy}{dx}=\dfrac{x}{y}$, with $y(0)=3$.
\tcblower
\textbf{Step 1 — General solution.} Separating: $y\,dy = x\,dx$. Integrating:
\[
\int y\,dy=\int x\,dx \;\Longrightarrow\; \frac{y^{2}}{2}=\frac{x^{2}}{2}+C \;\Longrightarrow\; y^{2}=x^{2}+2C.
\]
\textbf{Step 2 — Plug in the initial values.} With $x=0$, $y=3$: $9 = 0 + 2C$, so $C=\dfrac{9}{2}$.
\textbf{Step 3 — Write the solution.} $y^{2}=x^{2}+9$, hence $y=\pm\sqrt{x^{2}+9}$.
\textbf{Step 4 — Choose the branch.} Since $y(0)=3>0$, we must take the \emph{positive} branch:
\[
\boxed{\,y=\sqrt{x^{2}+9}\,},\qquad x\in\mathbb{R}.
\]
Check: $y(0)=\sqrt{9}=3$ \checkmark\ and $\dfrac{dy}{dx}=\dfrac{x}{\sqrt{x^{2}+9}}=\dfrac{x}{y}$ \checkmark.
\end{exoresolu}

\begin{attention}[The wrong branch trap]
When the general solution gives $y^{2}=\ldots$ or $|y|=\ldots$, the constant $C$ alone does \emph{not} fix the sign. The value of $C$ found may correspond to a branch that does not satisfy the initial condition. \textbf{Always verify by substitution}: compute $y(x_0)$ from your final answer and check it equals $y_0$, and check the sign of $y$ matches the data on the whole domain.
\end{attention}

\begin{attention}[Domain of the particular solution]
The particular solution must respect the domain: no division by zero, no logarithm of a non-positive number, and the correct branch of any square root. For example, if the general solution is $y=\dfrac{1}{C-x}$ and $y(0)=2$, then $C=\dfrac{1}{2}$ and the solution $y=\dfrac{1}{\tfrac12 - x}$ is only valid on $\left(-\infty,\tfrac12\right)$, the interval containing $x_0=0$.
\end{attention}

\courssub{Application 1: Exponential Growth and Decay}

Many quantities change at a rate proportional to their current size: a village population (more people, more births), a radioactive sample (more nuclei, more decays per second). This gives the most famous separable equation of all:
\[
\frac{dN}{dt}=kN.
\]

\begin{propriete}[Solution of the exponential equation]
The general solution of $\dfrac{dN}{dt}=kN$ is $N(t)=Ae^{kt}$, where $A=N(0)$ is the initial amount. If $k>0$ we have \cle{exponential growth}; if $k<0$, \cle{exponential decay}. (Derivation examinable: separate $\dfrac{dN}{dt}=kN$, integrate, exponentiate.)
\end{propriete}

\begin{exemplebox}[Population of a village]
A village near Bafoussam has $2\,000$ inhabitants and grows at a rate proportional to its population, with $k=0.03$ per year. Then $N(t)=2\,000\,e^{0.03t}$. After $10$ years: $N(10)=2\,000\,e^{0.3}\approx 2\,700$ inhabitants. The \emph{doubling time} satisfies $e^{0.03t}=2$, i.e. $t=\dfrac{\ln 2}{0.03}\approx 23.1$ years.
\end{exemplebox}

\begin{exemplebox}[Radioactive decay and half-life]
For a radioactive sample, $\dfrac{dN}{dt}=-\lambda N$ ($\lambda>0$), so $N(t)=N_0 e^{-\lambda t}$. The \cle{half-life} $T_{1/2}$ is the time for half the nuclei to decay:
\[
\frac{N_0}{2}=N_0 e^{-\lambda T_{1/2}} \;\Longrightarrow\; \boxed{\,T_{1/2}=\frac{\ln 2}{\lambda}\,}.
\]
Note the half-life does \emph{not} depend on $N_0$: this is the signature of exponential decay.
\end{exemplebox}

\courssub{Application 2: Newton's Law of Cooling}

\begin{propriete}[Newton's law of cooling]
The rate of change of the temperature $T$ of a body is proportional to the difference between $T$ and the (constant) ambient temperature $T_{\text{env}}$:
\[
\frac{dT}{dt}=-k\bigl(T-T_{\text{env}}\bigr),\qquad k>0.
\]
(Physical law stated; the derivation of its solution below is examinable.)
\end{propriete}

\textbf{Derivation.} Let $u=T-T_{\text{env}}$. Since $T_{\text{env}}$ is constant, $\dfrac{du}{dt}=\dfrac{dT}{dt}=-ku$. Hence $u=Ae^{-kt}$, i.e.
\[
T(t)=T_{\text{env}}+\bigl(T_0-T_{\text{env}}\bigr)e^{-kt},
\]
where $T_0=T(0)$. As $t\to\infty$, $e^{-kt}\to 0$, so $T\to T_{\text{env}}$: the body tends to the ambient temperature, exactly as experience tells us.

\begin{exoresolu}[A pot of ndolé cooling in the kitchen]
A pot of ndolé is at $95\ ^{\circ}\mathrm{C}$ when removed from the fire in a kitchen at $25\ ^{\circ}\mathrm{C}$. After $10$ minutes it is at $60\ ^{\circ}\mathrm{C}$. (a) Find $k$. (b) When will the ndolé reach $40\ ^{\circ}\mathrm{C}$?
\tcblower
The model gives $T(t)=25+70e^{-kt}$.
\textbf{(a)} At $t=10$: $60=25+70e^{-10k}$, so $e^{-10k}=\dfrac{35}{70}=\dfrac12$, hence
\[
k=\frac{\ln 2}{10}\approx 0.0693\ \text{min}^{-1}.
\]
\textbf{(b)} Set $T=40$: $40=25+70e^{-kt}\Rightarrow e^{-kt}=\dfrac{15}{70}=\dfrac{3}{14}$, so
\[
t=\frac{\ln(14/3)}{k}=\frac{10\ln(14/3)}{\ln 2}\approx 22.2\ \text{minutes}.
\]
Interpretation: the temperature drops fast at first (large difference with the kitchen), then ever more slowly, approaching $25\ ^{\circ}\mathrm{C}$ but never quite reaching it in finite time.
\end{exoresolu}

\begin{popfigure}
\begin{center}
\begin{tikzpicture}
\begin{axis}[
    width=11cm, height=6.5cm,
    xlabel={$t$ (min)}, ylabel={$T$ ($^{\circ}$C)},
    xmin=0, xmax=45, ymin=20, ymax=100,
    grid=both, grid style={popline, very thin},
    axis lines=left,
    legend style={at={(0.97,0.97)}, anchor=north east, font=\small},
    samples=100, domain=0:45,
]
\addplot[very thick, poporange] {25 + 70*exp(-0.0693*x)};
\addlegendentry{$T(t)=25+70e^{-0.0693t}$}
\addlegendentry{$T_{\text{env}}=25$}
\addplot[only marks, mark=*, popdark] coordinates {(0,95) (10,60) (22.2,40)};
\end{axis}
\end{tikzpicture}
\end{center}
\legende{Cooling curve of the pot of ndolé: exponential decay towards the ambient temperature $25\ ^{\circ}\mathrm{C}$ (horizontal asymptote).}
\end{popfigure}

\courssub{Application 3: Mixing Problems}

\begin{experience}[Salt in a tank]
A tank contains $100$ litres of brine with $500$ g of dissolved salt. Fresh water flows in at $5$ L/min and the well-stirred mixture flows out at the same rate of $5$ L/min (so the volume stays at $100$ L/min). Let $Q(t)$ be the mass of salt (in grams) at time $t$ (in minutes). The concentration in the tank is $\dfrac{Q}{100}$ g/L, so salt leaves at $5\times\dfrac{Q}{100}$ g/min, and none enters:
\[
\frac{dQ}{dt}=-\frac{Q}{20},\qquad Q(0)=500.
\]
Separating: $\dfrac{dQ}{Q}=-\dfrac{dt}{20}$, so $\ln|Q|=-\dfrac{t}{20}+C$, i.e. $Q=Ae^{-t/20}$. The initial condition gives $A=500$:
\[
Q(t)=500\,e^{-t/20}.
\]
For instance, the salt is reduced to $100$ g when $e^{-t/20}=\tfrac15$, i.e. $t=20\ln 5\approx 32.2$ min. If instead the incoming water contained salt at concentration $c$ g/L, the equation would become $\dfrac{dQ}{dt}=5c-\dfrac{Q}{20}$, still separable: $\dfrac{dQ}{5c-Q/20}=dt$.
\end{experience}

\begin{attention}[Keep units consistent]
In applications, keep units consistent throughout: temperature in $^{\circ}\mathrm{C}$, time in minutes (or seconds — but not both!), mass in grams, volume in litres. The constant $k$ then carries the correct units (e.g. min$^{-1}$), and mixing units is the most common source of dimensional errors in GCE modelling questions.
\end{attention}

\begin{attention}[Does the equation match the assumptions?]
When modelling, ensure the differential equation reflects the physical assumptions: Newton's law of cooling assumes a \emph{constant} ambient temperature; the mixing model assumes the tank is \emph{well stirred} and that inflow equals outflow (constant volume). If the assumptions change, the equation — and its solution — must change too.
\end{attention}

\courssub{Application 4 (GCE-style extension): Charging a Capacitor}

\begin{experience}[RC circuit]
A resistor $R$ and an initially uncharged capacitor $C$ are connected in series to a battery of e.m.f. $E$ at $t=0$. Kirchhoff's voltage law gives $E=Ri+\dfrac{q}{C}$ with $i=\dfrac{dq}{dt}$, hence
\[
R\frac{dq}{dt}+\frac{q}{C}=E
\quad\Longleftrightarrow\quad
\frac{dq}{dt}=\frac{CE-q}{RC},\qquad q(0)=0.
\]
This is separable: $\dfrac{dq}{CE-q}=\dfrac{dt}{RC}$. Integrating, $-\ln|CE-q|=\dfrac{t}{RC}+C_1$, so $CE-q=Ae^{-t/(RC)}$. With $q(0)=0$: $A=CE$, giving
\[
\boxed{\,q(t)=CE\Bigl(1-e^{-t/(RC)}\Bigr)\,}.
\]
As $t\to\infty$, $q\to CE$: the capacitor charges up to its limiting charge, with \cle{time constant} $\tau=RC$.
\end{experience}

\courssub{Logistic Growth: A Richer Separable Model}

Pure exponential growth cannot continue forever: a cassava crop is limited by space, water and nutrients. The \cle{logistic model} assumes the growth rate is proportional both to the population $P$ and to the ``room left'' $K-P$, where $K$ is the \emph{carrying capacity}:
\[
\frac{dP}{dt}=rP\left(1-\frac{P}{K}\right),\qquad P(0)=P_0.
\]
This is separable (partial fractions on $\dfrac{1}{P(K-P)}$), and its solution is
\[
P(t)=\frac{K}{1+\left(\dfrac{K-P_0}{P_0}\right)e^{-rt}}.
\]
As $t\to\infty$, $P\to K$: the population stabilises at the carrying capacity, producing the famous S-shaped curve.

\begin{aretenir}[Key points of this section]
\begin{itemize}
\item A \cle{particular solution} is obtained by substituting the initial condition $y(x_0)=y_0$ into the general solution and solving for $C$.
\item Always check the \emph{branch} (sign) and the \emph{domain} of the final answer by substituting back.
\item Exponential model: $N=N_0e^{kt}$; half-life $T_{1/2}=\dfrac{\ln 2}{\lambda}$.
\item Newton's cooling: $T=T_{\text{env}}+(T_0-T_{\text{env}})e^{-kt}$; limiting value $T_{\text{env}}$.
\item Mixing and RC circuits give equations of the form $\dfrac{dQ}{dt}=a-bQ$, solved by separation; limiting value $\dfrac{a}{b}$.
\item Logistic model: growth limited by a carrying capacity $K$.
\end{itemize}
\end{aretenir}

\begin{exercice}[Particular solutions and modelling]
\begin{enumerate}
\item Solve $\dfrac{dy}{dx}=2xy^{2}$ with $y(0)=1$, stating the largest interval of validity.
\item A radioactive sample loses $20\%$ of its nuclei in $5$ days. Find its half-life.
\item Water at $90\ ^{\circ}\mathrm{C}$ cools in a room at $20\ ^{\circ}\mathrm{C}$ with $k=0.05$ min$^{-1}$. After how long is it at $30\ ^{\circ}\mathrm{C}$?
\item A tank of $200$ L contains $1$ kg of salt. Brine at $0.5$ kg/L enters at $4$ L/min and the mixture leaves at $4$ L/min. Set up and solve the equation for the mass $Q(t)$ of salt, and find the limiting mass.
\end{enumerate}
\end{exercice}

\begin{corrige}[Exercise 1]
\textbf{1.} Separating: $\dfrac{dy}{y^{2}}=2x\,dx$, so $-\dfrac{1}{y}=x^{2}+C$. With $y(0)=1$: $C=-1$, hence $y=\dfrac{1}{1-x^{2}}$, valid on $(-1,1)$ (the interval containing $0$ where the denominator does not vanish).
\end{corrige}

\begin{savaistu}[Did you know?]
The same equation $\dfrac{dN}{dt}=kN$ governs the decay of carbon-14 used by archaeologists to date ancient remains, the growth of savings under compound interest, and the early spread of an epidemic. Mastering one separable equation unlocks dozens of real-world phenomena — a beautiful illustration of the unity of mathematics!
\end{savaistu}

\newpage

\coursec{Integration activity}

\coursec{Integration Activity: Cooling of a Pot of Ndolé}

\courssub{Aims of the Activity}
\begin{aretenir}[Learning Objectives]
\begin{itemize}
\item Recognise a first‑order differential equation that models a real‑world phenomenon.
\item Identify the equation as separable and perform the separation of variables.
\item Integrate to obtain the general solution, then apply an initial condition to find the particular solution.
\item Estimate the constant of proportionality from a second data point.
\item Use the model to make a prediction and critically discuss its limitations.
\end{itemize}
\end{aretenir}
This integration activity brings together all the skills developed in Lessons 61–66: meaning of a differential equation, recognition of separable form, separation of variables, integration, particular solutions, and application. By working on the cooling of a pot of \cle{ndolé} — a classic Cameroonian dish — you will see how mathematics describes everyday life in our kitchen.

\courssub{Situation}
\begin{experience}[Cooling of a Pot of Ndolé]
In a typical household in Yaoundé, a pot of ndolé has just been taken off the fire. Its initial temperature is \SI{95}{\ensuremath{{}^{\circ}\mathrm{C}}}. The kitchen (the environment) is at a constant temperature of \SI{25}{\ensuremath{{}^{\circ}\mathrm{C}}}. Experience tells us that the pot cools faster when it is very hot and more slowly as it approaches room temperature. This behaviour is modelled by \cle{Newton's Law of Cooling}:
\[
\frac{dT}{dt} = -k\,(T - T_{\text{env}}),
\]
where
\begin{itemize}
\item $T(t)$ is the temperature of the ndolé (in \si{\ensuremath{{}^{\circ}\mathrm{C}}}) at time $t$ (in minutes),
\item $T_{\text{env}} = 25$ is the constant ambient temperature,
\item $k > 0$ is a constant that depends on the pot, the liquid, and the surrounding air.
\end{itemize}
A family member measures the temperature after \SI{10}{\minute} and finds it to be \SI{70}{\ensuremath{{}^{\circ}\mathrm{C}}}. We assume that Newton's law holds exactly for this situation.
\end{experience}

\courssub{Tasks}
Work in groups of three or four. Answer the following questions in order, showing all steps clearly.

\begin{enumerate}
\item \textbf{Write the differential equation} for $T(t)$ using the given data. Identify the type of equation (order, linearity, separability).
\item \textbf{Separate the variables} $T$ and $t$ and integrate both sides to obtain the \emph{general solution}. Express the result in the form $T(t) = T_{\text{env}} + C e^{-kt}$, where $C$ is an arbitrary constant.
\item \textbf{Apply the initial condition} $T(0) = 95$ to determine the constant $C$ and write the \emph{particular solution} that satisfies this condition.
\item \textbf{Estimate the constant $k$} using the second measurement $T(10) = 70$. Give the exact expression for $k$ (in terms of a logarithm) and a numerical approximation to three significant figures.
\item \textbf{Predict the temperature} after \SI{30}{\minute} using your model. Round your answer to one decimal place.
\item \textbf{Discuss the limitations} of the model. Consider at least three factors that could make the real cooling curve deviate from the exponential prediction.
\end{enumerate}

\courssub{Model Answer}
\begin{exemplebox}[Complete Worked Solution]
\textbf{Task 1 – Differential Equation}\\
The situation gives $T_{\text{env}} = 25$. Substituting into Newton's law:
\[
\frac{dT}{dt} = -k\,(T - 25).
\]
This is a \emph{first‑order} ordinary differential equation (only the first derivative appears). It is \emph{linear} in $T$ (can be written as $\frac{dT}{dt} + kT = 25k$) and, crucially, it is \emph{separable} because the right‑hand side is a product of a function of $T$ and a function of $t$ (the latter being the constant $-k$).

\textbf{Task 2 – Separation and General Solution}\\
Separate the variables:
\[
\frac{dT}{T - 25} = -k\,dt.
\]
Integrate both sides:
\[
\int \frac{dT}{T - 25} = \int -k\,dt \quad\Longrightarrow\quad \ln|T - 25| = -kt + C_1,
\]
where $C_1$ is an arbitrary constant. Exponentiate:
\[
|T - 25| = e^{C_1} e^{-kt}.
\]
Since the pot cools from $95^\circ\mathrm{C}$ towards $25^\circ\mathrm{C}$, we have $T > 25$ for all $t \ge 0$; hence $T-25 > 0$ and the absolute value can be dropped. Let $C = e^{C_1}$ (so $C > 0$ in this physical context, but the algebraic form allows any non‑zero constant). Thus
\[
T - 25 = C e^{-kt} \quad\Longrightarrow\quad \boxed{T(t) = 25 + C e^{-kt}}.
\]

\textbf{Task 3 – Particular Solution from $T(0)=95$}\\
At $t=0$, $T=95$:
\[
95 = 25 + C e^{0} = 25 + C \quad\Longrightarrow\quad C = 70.
\]
Hence the particular solution is
\[
\boxed{T(t) = 25 + 70 e^{-kt}}.
\]

\textbf{Task 4 – Estimating $k$ from $T(10)=70$}\\
Substitute $t=10$, $T=70$:
\[
70 = 25 + 70 e^{-10k} \;\Longrightarrow\; 45 = 70 e^{-10k} \;\Longrightarrow\; e^{-10k} = \frac{45}{70} = \frac{9}{14}.
\]
Take natural logarithms:
\[
-10k = \ln\!\left(\frac{9}{14}\right) \;\Longrightarrow\; k = -\frac{1}{10}\ln\!\left(\frac{9}{14}\right) = \frac{1}{10}\ln\!\left(\frac{14}{9}\right).
\]
Exact form: $\displaystyle k = \frac{1}{10}\ln\!\left(\frac{14}{9}\right)\ \mathrm{min}^{-1}$.
Numerical approximation:
\[
\ln\!\left(\frac{14}{9}\right) \approx 0.4559 \quad\Longrightarrow\quad k \approx 0.0456\ \mathrm{min}^{-1} \quad\text{(three significant figures)}.
\]

\textbf{Task 5 – Temperature after 30 Minutes}\\
Use $T(t) = 25 + 70 e^{-kt}$ with the exact $k$:
\[
T(30) = 25 + 70 e^{-30k} = 25 + 70 \left(e^{-10k}\right)^3 = 25 + 70 \left(\frac{9}{14}\right)^3.
\]
Compute $\left(\frac{9}{14}\right)^3 = \frac{729}{2744} \approx 0.2657$.
\[
T(30) \approx 25 + 70 \times 0.2657 = 25 + 18.60 = 43.6^\circ\mathrm{C}.
\]
\boxed{T(30) $\approx$ 43.6^\circ\mathrm{C}}.

\textbf{Task 6 – Limitations of the Model}\\
\begin{enumerate}
\item \textbf{Constant ambient temperature:} In a real kitchen, opening doors, other cooking, or the pot itself warming the air can change $T_{\text{env}}$.
\item \textbf{Uniform temperature assumption:} The model assumes the whole pot is at a single temperature $T$. In reality, there are temperature gradients (hotter at the bottom, cooler at the surface), especially if the ndolé is not stirred.
\item \textbf{Constant $k$:} The cooling constant $k$ depends on the heat transfer coefficient, which varies with the temperature difference (natural convection changes), the exposed surface area (evaporation reduces liquid level), and whether a lid is used.
\item \textbf{Phase change and evaporation:} As the ndolé cools, water evaporates, carrying away latent heat. This adds an extra cooling mechanism not captured by the simple linear law.
\item \textbf{Measurement errors:} The thermometer has finite precision, and the 10‑minute reading may not be exact.
\end{enumerate}
These factors mean the exponential model is a \emph{first approximation}; a more sophisticated model would include variable $k$, heat conduction within the pot, and mass loss.
\end{exemplebox}

\begin{popfigure}
\begin{tikzpicture}
\begin{axis}[
    width=\linewidth,
    height=8cm,
    xlabel={Time $t$ (min)},
    ylabel={Temperature $T$ ($^\circ\mathrm{C}$)},
    xmin=0, xmax=40,
    ymin=20, ymax=100,
    grid=major,
    grid style={popline},
    tick label style={font=\small},
    label style={font=\small},
    legend style={at={(0.98,0.98)}, anchor=north east, font=\small, fill=white, draw=popline}
]
\addplot[popblue, thick, domain=0:40, samples=200] {25 + 70*exp(-0.0456*x)};
\addlegendentry{Model $T(t)=25+70e^{-0.0456t}$}
\addplot[only marks, mark=*, poporange] coordinates {(0,95) (10,70) (30,43.6)};
\addlegendentry{Data / Prediction}
\draw[dashed, popmuted] (0,25) -- (40,25) node[pos=0.9, above, font=\tiny] {$T_{\text{env}}=25^\circ\mathrm{C}$};
\end{axis}
\end{tikzpicture}
\legende{Cooling curve of the ndolé pot. The model (blue) passes through the initial point $(0,95)$ and the 10‑minute measurement $(10,70)$. The predicted temperature at 30 minutes is $(30,43.6)$. The horizontal dashed line is the ambient temperature.}
\end{popfigure}

\begin{attention}[Common Pitfalls]
\begin{itemize}
\item Forgetting that $T-25>0$ when removing the absolute value after integration.
\item Mixing up the sign of $k$; the law has a minus sign, so $k$ must be positive.
\item Using $\log_{10}$ instead of $\ln$ when solving for $k$; the integration gives natural logarithms.
\item Rounding $k$ too early, which propagates error into the 30‑minute prediction.
\item Stating limitations that are not physical (e.g., ``the pot might be stolen'') instead of modelling assumptions.
\end{itemize}
\end{attention}

\begin{savaistu}[Did You Know?]
Newton's law of cooling was first formulated by Sir Isaac Newton in 1701. In Cameroon, it is used not only in cooking but also in forensic science to estimate the time of death (the ``body cooling'' method) and in engineering to design cooling systems for transformers in the ENEO power grid. The same mathematics describes the discharge of a capacitor through a resistor (RC circuit) — a perfect example of the unity of science!
\end{savaistu}

\newpage

\coursec{Methods and worked examples}

\courssub{Method 1: Verifying that a Function is a Solution of a Differential Equation}

\begin{methode}[Verifying a proposed solution]
To check whether a given function $y = g(x)$ is a solution of the first order differential equation $\dfrac{dy}{dx} = f(x, y)$ on an interval $I$:
\begin{enumerate}
\item Differentiate the proposed function to obtain $\dfrac{dy}{dx} = g'(x)$, stating the interval on which $g$ is differentiable.
\item Compute the right-hand side $f\big(x, g(x)\big)$ by substituting $y = g(x)$ into the equation.
\item Compare the two expressions. If $g'(x) = f\big(x, g(x)\big)$ for \emph{every} $x$ in $I$, then $g$ is a solution on $I$; otherwise it is not.
\item If an initial condition $y(x_0) = y_0$ is given, check separately that $g(x_0) = y_0$.
\end{enumerate}
\textbf{Reflex:} never substitute only at one point --- the identity must hold for all $x$ in the interval.
\end{methode}

\begin{exoresolu}[Verifying solutions of $y' = 2xy$]
Show that for every constant $C \in \mathbb{R}$, the function $y = Ce^{x^2}$ is a solution of the differential equation $\dfrac{dy}{dx} = 2xy$ on $\mathbb{R}$. Determine which of these solutions satisfies $y(0) = 3$.
\tcblower
\textbf{Step 1: Differentiate.} The function $y = Ce^{x^2}$ is differentiable on $\mathbb{R}$ and, by the chain rule,
\[ \frac{dy}{dx} = C \cdot 2x\, e^{x^2} = 2Cx e^{x^2}. \]
\textbf{Step 2: Compute the right-hand side.} Substituting $y = Ce^{x^2}$ into $2xy$:
\[ 2xy = 2x \cdot Ce^{x^2} = 2Cx e^{x^2}. \]
\textbf{Step 3: Compare.} For every $x \in \mathbb{R}$,
\[ \frac{dy}{dx} = 2Cx e^{x^2} = 2xy. \]
The identity holds for all $x \in \mathbb{R}$ and for every value of $C$, so $y = Ce^{x^2}$ is a solution on $\mathbb{R}$ for every $C \in \mathbb{R}$. This family is in fact the \cle{general solution} of the equation.

\textbf{Step 4: Initial condition.} We require $y(0) = 3$:
\[ Ce^{0} = 3 \quad\Longrightarrow\quad C = 3. \]
Hence the \cle{particular solution} satisfying $y(0) = 3$ is $\boxed{y = 3e^{x^2}}$.
\end{exoresolu}

\courssub{Method 2: Recognising a Separable Equation}

\begin{methode}[Testing for separability]
A first order equation $\dfrac{dy}{dx} = f(x, y)$ is \cle{separable} if $f(x, y)$ can be written as a product (or quotient) of a function of $x$ alone and a function of $y$ alone:
\[ \frac{dy}{dx} = g(x)\,h(y). \]
\begin{enumerate}
\item Write the equation in the form $\dfrac{dy}{dx} = f(x, y)$.
\item Try to factorise $f(x, y)$: look for common factors, use identities (e.g.\ $x^2 + x^2 y = x^2(1+y)$), or exponentials ($e^{x+y} = e^x e^y$).
\item If $f(x,y) = g(x)h(y)$, the equation is separable. If the variables cannot be disentangled (e.g.\ $\dfrac{dy}{dx} = x + y$ or $\dfrac{dy}{dx} = \sin(xy)$), it is \textbf{not} separable.
\item Note the special cases: $\dfrac{dy}{dx} = g(x)$ (here $h(y) = 1$) and $\dfrac{dy}{dx} = h(y)$ (here $g(x) = 1$) are both separable.
\end{enumerate}
\end{methode}

\begin{exoresolu}[Which equations are separable?]
For each equation, state whether it is separable; if it is, write it in the form $\dfrac{dy}{dx} = g(x)h(y)$.
\begin{enumerate}
\item[(a)] $\dfrac{dy}{dx} = \dfrac{x^2}{y}$ \qquad (b) $\dfrac{dy}{dx} = e^{x - y}$ \qquad (c) $\dfrac{dy}{dx} = xy + x$ \qquad (d) $\dfrac{dy}{dx} = x + y^2$
\end{enumerate}
\tcblower
\textbf{(a)} $\dfrac{dy}{dx} = \dfrac{x^2}{y} = x^2 \cdot \dfrac{1}{y}$. This is $g(x)h(y)$ with $g(x) = x^2$ and $h(y) = \dfrac{1}{y}$. \textbf{Separable.}

\textbf{(b)} Using the law of indices, $e^{x-y} = e^x \cdot e^{-y}$. Hence $\dfrac{dy}{dx} = e^x \cdot e^{-y}$, with $g(x) = e^x$, $h(y) = e^{-y}$. \textbf{Separable.}

\textbf{(c)} Factorise the right-hand side: $xy + x = x(y + 1)$. Hence $\dfrac{dy}{dx} = x(y+1)$, with $g(x) = x$, $h(y) = y+1$. \textbf{Separable.}

\textbf{(d)} $x + y^2$ is a \emph{sum} mixing $x$ and $y$; it cannot be written as a product $g(x)h(y)$. (Indeed, if $x + y^2 = g(x)h(y)$ for all $x, y$, fixing $y$ and varying $x$ would force $h$ to be constant, which is impossible.) \textbf{Not separable.}

\textbf{Reflex:} factorisation is the key skill --- a sum can often be factored into a product, as in (c), so do not reject an equation before trying to factorise.
\end{exoresolu}

\courssub{Method 3: Solving a Separable Equation --- General Solution}

\begin{methode}[Separation of variables]
To solve $\dfrac{dy}{dx} = g(x)\,h(y)$:
\begin{enumerate}
\item \textbf{Separate:} divide by $h(y)$ (noting where $h(y) = 0$) and write
\[ \frac{1}{h(y)}\, dy = g(x)\, dx. \]
All $y$'s with $dy$ on one side, all $x$'s with $dx$ on the other.
\item \textbf{Integrate both sides:}
\[ \int \frac{1}{h(y)}\, dy = \int g(x)\, dx. \]
\item \textbf{One constant only:} put a constant of integration on \emph{one} side only (the two constants merge into one).
\item \textbf{Solve for $y$} if possible, to obtain the general solution explicitly; otherwise leave it in implicit form.
\item \textbf{Check lost solutions:} constant solutions $y = y_0$ with $h(y_0) = 0$ may have been excluded in step 1; check whether they satisfy the equation.
\end{enumerate}
\end{methode}

\begin{exoresolu}[General solution of $\dfrac{dy}{dx} = \dfrac{x}{y}$]
Find the general solution of the differential equation $\dfrac{dy}{dx} = \dfrac{x}{y}$, and describe the family of solution curves.
\tcblower
\textbf{Step 1: Separate.} The equation is separable with $g(x) = x$, $h(y) = \dfrac{1}{y}$. Multiplying both sides by $y$:
\[ y\, dy = x\, dx. \]
\textbf{Step 2: Integrate.}
\[ \int y\, dy = \int x\, dx \quad\Longrightarrow\quad \frac{y^2}{2} = \frac{x^2}{2} + C_1. \]
\textbf{Step 3: Simplify.} Multiplying by 2 and writing $C = 2C_1$:
\[ y^2 = x^2 + C, \qquad\text{i.e.}\qquad y^2 - x^2 = C. \]
\textbf{Step 4: Explicit form.} $y = \pm\sqrt{x^2 + C}$, the sign being fixed by the branch on which the solution lives.

\textbf{Interpretation:} the solution curves are the hyperbolas $y^2 - x^2 = C$ (rectangular hyperbolas with asymptotes $y = \pm x$), together with the lines $y = \pm x$ as limiting cases $C = 0$. Each initial condition $(x_0, y_0)$ with $y_0 \neq 0$ selects exactly one branch.

\textbf{Check:} differentiating $y^2 = x^2 + C$ implicitly gives $2y\dfrac{dy}{dx} = 2x$, i.e.\ $\dfrac{dy}{dx} = \dfrac{x}{y}$. \checkmark
\end{exoresolu}

\courssub{Method 4: Equations Involving Logarithms and Exponentials}

\begin{methode}[Handling $\int \frac{dy}{y}$ and removing logarithms]
When separation produces $\displaystyle\int \frac{dy}{y} = \ln|y|$, proceed as follows:
\begin{enumerate}
\item Write $\ln|y| = F(x) + C$ where $F(x) = \int g(x)\,dx$.
\item Exponentiate: $|y| = e^{F(x) + C} = e^{C}\, e^{F(x)}$.
\item Absorb the sign and the constant: since $e^C > 0$ and $\pm e^C$ ranges over $\mathbb{R}\setminus\{0\}$, write $y = A e^{F(x)}$ with $A \in \mathbb{R}\setminus\{0\}$; then check whether $y = 0$ is also a solution and include it by allowing $A = 0$ when valid.
\end{enumerate}
\textbf{Key formulas:} $e^{\ln u} = u$ ($u>0$); \; $e^{a+b} = e^a e^b$; \; $\ln|y|$ requires the absolute value.
\end{methode}

\begin{exoresolu}[General solution of $\dfrac{dy}{dx} = \dfrac{2y}{x}$ ($x > 0$)]
Find the general solution of $\dfrac{dy}{dx} = \dfrac{2y}{x}$ for $x > 0$, giving $y$ explicitly in terms of $x$.
\tcblower
\textbf{Step 1: Separate.} For $y \neq 0$, divide by $y$ and multiply by $dx$:
\[ \frac{dy}{y} = \frac{2\, dx}{x}. \]
\textbf{Step 2: Integrate.}
\[ \int \frac{dy}{y} = \int \frac{2}{x}\, dx \quad\Longrightarrow\quad \ln|y| = 2\ln x + C \quad (x > 0). \]
\textbf{Step 3: Exponentiate.}
\[ |y| = e^{2\ln x + C} = e^{C}\, e^{2\ln x} = e^{C}\, x^{2}. \]
Writing $A = \pm e^{C}$ (so $A \neq 0$):
\[ y = A x^{2}. \]
\textbf{Step 4: Lost solution.} In Step 1 we assumed $y \neq 0$. The constant function $y = 0$ gives $\dfrac{dy}{dx} = 0$ and $\dfrac{2y}{x} = 0$, so it \emph{is} a solution; it corresponds to $A = 0$.

\textbf{Conclusion:} the general solution is $\boxed{y = Ax^2,\ A \in \mathbb{R}}$, a family of parabolas through the origin (plus the $x$-axis).

\textbf{Check:} $y = Ax^2 \Rightarrow \dfrac{dy}{dx} = 2Ax = \dfrac{2(Ax^2)}{x} = \dfrac{2y}{x}$. \checkmark
\end{exoresolu}

\courssub{Method 5: Finding Particular Solutions from Initial Conditions}

\begin{methode}[Initial value problems]
To solve the initial value problem $\dfrac{dy}{dx} = g(x)h(y)$, $y(x_0) = y_0$:
\begin{enumerate}
\item Find the \textbf{general solution} first (separate, integrate, one constant $C$).
\item \textbf{Substitute} $x = x_0$ and $y = y_0$ into the general solution to determine $C$.
\item Write the \textbf{particular solution}, stating its interval of validity (the largest interval containing $x_0$ on which the solution is defined and satisfies the equation).
\item \textbf{Verify} by differentiating and by re-checking the initial condition.
\end{enumerate}
\begin{attention}[Common traps]
\begin{itemize}
\item Substituting the initial condition \emph{before} integrating --- the constant can only be found \emph{after} integration.
\item Forgetting the absolute value in $\ln|y|$ and then losing sign information when the initial value is negative.
\item Giving $y = \pm\sqrt{\cdots}$ without using the initial condition to choose the sign.
\end{itemize}
\end{attention}
\end{methode}

\begin{exoresolu}[Solving $\dfrac{dy}{dx} = \dfrac{x^2}{y^2}$ with $y(0) = 2$]
Solve the initial value problem $\dfrac{dy}{dx} = \dfrac{x^2}{y^2}$, $y(0) = 2$, giving $y$ explicitly as a function of $x$.
\tcblower
\textbf{Step 1: Separate and integrate.}
\[ y^2\, dy = x^2\, dx \quad\Longrightarrow\quad \int y^2\, dy = \int x^2\, dx \quad\Longrightarrow\quad \frac{y^3}{3} = \frac{x^3}{3} + C_1. \]
Multiplying by 3 and setting $C = 3C_1$: $\;y^3 = x^3 + C$.

\textbf{Step 2: Apply the initial condition.} Put $x = 0$, $y = 2$:
\[ 2^3 = 0^3 + C \quad\Longrightarrow\quad C = 8. \]
\textbf{Step 3: Particular solution.} $y^3 = x^3 + 8$, and since the cube root is unique,
\[ \boxed{y = \sqrt[3]{x^3 + 8}}. \]
This is defined for all $x \in \mathbb{R}$, and $y \neq 0$ except at $x = -2$; on any interval not containing $-2$ the equation is satisfied.

\textbf{Step 4: Verify.} Differentiating $y^3 = x^3 + 8$ implicitly: $3y^2 \dfrac{dy}{dx} = 3x^2$, so $\dfrac{dy}{dx} = \dfrac{x^2}{y^2}$. \checkmark\ Also $y(0) = \sqrt[3]{8} = 2$. \checkmark
\end{exoresolu}

\begin{exoresolu}[A trigonometric equation: $\dfrac{dy}{dx} = \dfrac{\cos x}{y}$, $y(0) = -3$]
Solve $\dfrac{dy}{dx} = \dfrac{\cos x}{y}$ given that $y = -3$ when $x = 0$. Express $y$ explicitly in terms of $x$ and state the largest interval containing $0$ on which the solution is valid.
\tcblower
\textbf{Step 1: Separate and integrate.}
\[ y\, dy = \cos x\, dx \quad\Longrightarrow\quad \frac{y^2}{2} = \sin x + C_1 \quad\Longrightarrow\quad y^2 = 2\sin x + C. \]
\textbf{Step 2: Initial condition.} $x = 0$, $y = -3$:
\[ (-3)^2 = 2\sin 0 + C \quad\Longrightarrow\quad C = 9. \]
Hence $y^2 = 2\sin x + 9$.
\textbf{Step 3: Choose the sign.} Since $y(0) = -3 < 0$, we take the \emph{negative} square root:
\[ \boxed{y = -\sqrt{2\sin x + 9}}. \]
\textbf{Step 4: Interval of validity.} Since $-1 \le \sin x \le 1$, we have $2\sin x + 9 \ge 7 > 0$ for all $x$, so the expression under the root never vanishes: the solution is valid on the whole of $\mathbb{R}$.

\textbf{Verify:} $2y y' = 2\cos x \Rightarrow y' = \dfrac{\cos x}{y}$. \checkmark\; $y(0) = -\sqrt{9} = -3$. \checkmark

\textbf{Examiner's remark:} choosing $y = +\sqrt{2\sin x + 9}$ would contradict the initial condition; the sign is decided by the data, not by preference.
\end{exoresolu}

\courssub{Method 6: Modelling with Separable Equations (Applications)}

\begin{methode}[Setting up and solving applied problems]
\begin{enumerate}
\item \textbf{Translate} the wording into a differential equation: "rate of change proportional to ..." means $\dfrac{dy}{dt} = ky$ (growth/decay); "proportional to the difference" means $\dfrac{dy}{dt} = k(y - A)$.
\item \textbf{Identify the variables and units}, and write down all given data as initial/boundary conditions.
\item \textbf{Separate and integrate} to get the general solution.
\item Use the conditions to find the constant of integration $C$ \emph{and} the proportionality constant $k$ (two conditions are needed in general).
\item \textbf{Answer the question asked} (a value, a time, a half-life...), with units, and check the answer is physically sensible.
\end{enumerate}
\textbf{Key model:} $\dfrac{dy}{dt} = ky \;\Longrightarrow\; y = y_0 e^{kt}$: exponential growth if $k>0$, decay if $k<0$.
\end{methode}

\begin{exoresolu}[Population growth of bacteria]
A bacterial culture in a laboratory in Yaounde grows at a rate proportional to the number $N$ of bacteria present. Initially there are $1000$ bacteria, and after $2$ hours there are $3000$.
\begin{enumerate}
\item[(a)] Show that $N = 1000\, e^{kt}$ where $k = \tfrac{1}{2}\ln 3$.
\item[(b)] Find the number of bacteria after $6$ hours.
\item[(c)] Find the time at which the population reaches $27000$.
\end{enumerate}
\tcblower
\textbf{(a) Setting up.} "Rate of growth proportional to the population" gives
\[ \frac{dN}{dt} = kN. \]
Separating (for $N > 0$): $\dfrac{dN}{N} = k\, dt$, and integrating:
\[ \ln N = kt + C \quad\Longrightarrow\quad N = Ae^{kt}. \]
Initial condition $N(0) = 1000$: $A = 1000$, so $N = 1000e^{kt}$.

Condition $N(2) = 3000$: $3000 = 1000e^{2k} \Rightarrow e^{2k} = 3 \Rightarrow 2k = \ln 3$, hence $k = \tfrac{1}{2}\ln 3 \approx 0.549\ \mathrm{h^{-1}}$. Therefore
\[ N = 1000\, e^{kt} \quad\text{with}\quad k = \tfrac12 \ln 3. \qquad\blacksquare \]

\textbf{(b) After 6 hours:}
\[ N(6) = 1000 e^{6k} = 1000 e^{3\ln 3} = 1000 \times 3^3 = 27000 \text{ bacteria}. \]
(Note $e^{3\ln 3} = (e^{\ln 3})^3 = 27$.)

\textbf{(c) Time to reach 27000:} from (b) this is exactly $t = 6$ hours. More generally, solving $1000 e^{kt} = 27000$ gives $e^{kt} = 27 = 3^3 = e^{6k}$, so $t = 6$ h. \checkmark

\textbf{Sanity check:} the population triples every 2 hours ($1000 \to 3000 \to 9000 \to 27000$), consistent with $e^{2k} = 3$.
\end{exoresolu}

\begin{exoresolu}[Newton's law of cooling]
A pot of \emph{ndi\'e} (pepper soup) is at $90\,^\circ\mathrm{C}$ when removed from the fire in a kitchen kept at $30\,^\circ\mathrm{C}$. According to Newton's law of cooling, the temperature $\theta$ (in $^\circ$C) satisfies $\dfrac{d\theta}{dt} = -k(\theta - 30)$, where $t$ is in minutes and $k > 0$. After $5$ minutes the soup is at $70\,^\circ\mathrm{C}$.
\begin{enumerate}
\item[(a)] Show that $\theta = 30 + 60e^{-kt}$ and find the exact value of $k$.
\item[(b)] Find the temperature after $10$ minutes.
\end{enumerate}
\tcblower
\textbf{(a)} Let $u = \theta - 30$ (the temperature excess). Then $\dfrac{du}{dt} = \dfrac{d\theta}{dt} = -ku$, a standard decay equation. Separating:
\[ \frac{du}{u} = -k\, dt \quad\Longrightarrow\quad \ln|u| = -kt + C \quad\Longrightarrow\quad u = Ae^{-kt}. \]
Hence $\theta = 30 + Ae^{-kt}$. At $t = 0$, $\theta = 90$: $90 = 30 + A \Rightarrow A = 60$, so
\[ \theta = 30 + 60e^{-kt}. \qquad\blacksquare \]
At $t = 5$, $\theta = 70$: $\;70 = 30 + 60e^{-5k} \Rightarrow e^{-5k} = \dfrac{40}{60} = \dfrac{2}{3}$, so
\[ -5k = \ln\tfrac{2}{3} \quad\Longrightarrow\quad k = \frac{1}{5}\ln\frac{3}{2} \approx 0.0811\ \mathrm{min^{-1}}. \]

\textbf{(b) After 10 minutes:}
\[ \theta(10) = 30 + 60e^{-10k} = 30 + 60\left(e^{-5k}\right)^2 = 30 + 60\left(\tfrac{2}{3}\right)^2 = 30 + 60 \times \tfrac{4}{9} = 30 + \tfrac{80}{3} = \frac{170}{3} \approx 56.7\,^\circ\mathrm{C}. \]

\textbf{Sanity check:} the temperature decreases towards the room temperature $30\,^\circ\mathrm{C}$ as $t \to \infty$ (since $e^{-kt} \to 0$), which is physically correct; and $56.7 < 70$, consistent with continued cooling.
\end{exoresolu}

\begin{exoresolu}[GCE-style synthesis: $\dfrac{dy}{dx} = (y+1)\sin x$]
\begin{enumerate}
\item[(a)] Find the general solution of $\dfrac{dy}{dx} = (y+1)\sin x$.
\item[(b)] Find the particular solution for which $y = 0$ when $x = 0$, expressing $y$ explicitly in terms of $x$.
\item[(c)] Find the exact value of $y$ when $x = \dfrac{\pi}{2}$.
\end{enumerate}
\tcblower
\textbf{(a)} The equation is separable: $g(x) = \sin x$, $h(y) = y + 1$. For $y \neq -1$:
\[ \frac{dy}{y+1} = \sin x\, dx \quad\Longrightarrow\quad \int \frac{dy}{y+1} = \int \sin x\, dx \quad\Longrightarrow\quad \ln|y+1| = -\cos x + C. \]
Exponentiating: $|y+1| = e^{C}e^{-\cos x}$, so $y + 1 = Ae^{-\cos x}$ with $A \neq 0$; the constant solution $y = -1$ (excluded when dividing) corresponds to $A = 0$. General solution:
\[ y = Ae^{-\cos x} - 1, \qquad A \in \mathbb{R}. \]

\textbf{(b)} Put $x = 0$, $y = 0$: $\;0 = Ae^{-\cos 0} - 1 = Ae^{-1} - 1 \Rightarrow A = e$. Hence
\[ \boxed{y = e^{\,1 - \cos x} - 1}. \]

\textbf{(c)} At $x = \dfrac{\pi}{2}$, $\cos x = 0$:
\[ y = e^{1 - 0} - 1 = e - 1 \approx 1.718. \]

\textbf{Verify:} $\dfrac{dy}{dx} = e^{1-\cos x}\cdot \sin x = (y+1)\sin x$. \checkmark\; At $x=0$: $y = e^{0} - 1 = 0$. \checkmark
\end{exoresolu}

\newpage

\coursec{Exercises}

\courssub{I check my knowledge}

\begin{exercice}[MCQ: Recognising separable equations]
Which of the following first order differential equations is separable?
\begin{enumerate}
\item $\dfrac{dy}{dx} = x + y$
\item $\dfrac{dy}{dx} = \dfrac{x^2}{y}$
\item $\dfrac{dy}{dx} = \sin(x+y)$
\item $\dfrac{dy}{dx} = e^{x}y^2$
\end{enumerate}
\end{exercice}

\begin{exercice}[True/False with justification]
State whether the following statement is true or false. Justify your answer.
\\ ``The differential equation $\dfrac{dy}{dx} = \dfrac{x^2 + y^2}{xy}$ is separable.''
\end{exercice}

\begin{exercice}[Definition]
Define a \emph{first order separable differential equation}. Give the general form and explain the condition that allows the variables to be separated.
\end{exercice}

\begin{exercice}[Direct integration]
Solve the differential equation $\dfrac{dy}{dx} = \dfrac{2x}{y}$, expressing the general solution implicitly.
\end{exercice}

\courssub{I practise}

\begin{exercice}[Initial value problem]
Solve the initial value problem:
\[
\frac{dy}{dx} = \frac{3x^2}{2y},\qquad y(1)=2.
\]
Give the explicit particular solution.
\end{exercice}

\begin{exercice}[Exponential growth: bacterial culture]
A bacterial culture grows at a rate proportional to its size. If the population doubles in 3 hours, find the time needed for the population to triple.
\end{exercice}

\begin{exercice}[Mixing problem]
A tank initially contains 100 L of brine with 5 kg of dissolved salt. Pure water enters the tank at a rate of 5 L/min, and the well‑stirred mixture leaves at the same rate.
\begin{enumerate}
\item Set up the differential equation for the amount of salt $S(t)$ (in kg) in the tank at time $t$ minutes.
\item Solve the differential equation and find $S(t)$.
\end{enumerate}
\end{exercice}

\begin{exercice}[Newton's law of cooling]
The temperature $T$ (in $^\circ\mathrm{C}$) of a cup of coffee obeys the differential equation
\[
\frac{dT}{dt} = -0.07\,(T-20),
\]
where $t$ is in minutes. If the coffee is initially at $85^\circ\mathrm{C}$, how long will it take to cool to $50^\circ\mathrm{C}$?
\end{exercice}

\courssub{I prepare for the GCE}

\begin{exercice}[Recognition and substitution — 5 marks]
Determine whether the differential equation
\[
\frac{dy}{dx} = \frac{x^2 + y^2}{xy}
\]
is separable. If it is not, suggest a substitution that transforms it into a separable equation. (Do not solve the transformed equation.)
\end{exercice}

\begin{exercice}[General solution and constant solutions — 6 marks]
Find the general solution of the differential equation
\[
\frac{dy}{dx} = \frac{y\,\ln y}{x},\qquad x>0,\; y>0.
\]
Discuss any constant solutions that may have been lost during the separation process.
\end{exercice}

\begin{exercice}[Radioactive decay — 7 marks]
A radioactive substance has a half‑life of 12 years.
\begin{enumerate}
\item Write the differential equation that models the mass $m(t)$ of the substance at time $t$ years.
\item Solve the differential equation to obtain $m(t)$ in terms of the initial mass $m_0$.
\item Find the time required for the mass to reduce to 10\% of its original value.
\end{enumerate}
\end{exercice}

\newpage

\coursec{Solutions to the exercises}

\begin{corrige}[Exercise 1 — MCQ: Recognising separable equations]
A first order equation is separable if it can be written in the form $\dfrac{dy}{dx}=f(x)\,g(y)$, i.e. the right-hand side is a product of a function of $x$ alone and a function of $y$ alone.

\begin{itemize}
\item \textbf{(a)} $\dfrac{dy}{dx}=x+y$: the sum $x+y$ cannot be factored as $f(x)g(y)$. \textbf{Not separable.}
\item \textbf{(b)} $\dfrac{dy}{dx}=\dfrac{x^2}{y}=x^2\cdot\dfrac{1}{y}$: product of $f(x)=x^2$ and $g(y)=\dfrac{1}{y}$. \textbf{Separable.}
\item \textbf{(c)} $\dfrac{dy}{dx}=\sin(x+y)$: the argument $x+y$ mixes the variables inside the sine; no factorisation $f(x)g(y)$ exists. \textbf{Not separable.}
\item \textbf{(d)} $\dfrac{dy}{dx}=e^{x}y^{2}$: product of $f(x)=e^{x}$ and $g(y)=y^{2}$. \textbf{Separable.}
\end{itemize}

\textbf{Answer: (b) and (d).}
\end{corrige}

\begin{corrige}[Exercise 2 — True/False with justification]
\textbf{The statement is FALSE.}

\textbf{Justification.} Suppose, for contradiction, that
\[
\frac{x^{2}+y^{2}}{xy}=f(x)\,g(y)
\]
for some functions $f$ and $g$. Expanding,
\[
\frac{x^{2}+y^{2}}{xy}=\frac{x}{y}+\frac{y}{x}.
\]
This is a \emph{sum} of a term involving $\dfrac{x}{y}$ and a term involving $\dfrac{y}{x}$; the variables are coupled through the ratios $\dfrac{x}{y}$ and $\dfrac{y}{x}$ and cannot be split into a product of a function of $x$ alone by a function of $y$ alone. For instance, at $(x,y)=(1,1)$ the right-hand side equals $2$, at $(1,2)$ it equals $\dfrac{5}{2}$ and at $(2,1)$ it equals $\dfrac{5}{2}$; a product $f(x)g(y)$ satisfying $f(1)g(1)=2$, $f(1)g(2)=\tfrac52$, $f(2)g(1)=\tfrac52$ would force $f(2)g(2)=\dfrac{(5/2)^2}{2}=\dfrac{25}{8}$, whereas the actual value is $\dfrac{4+4}{2}=4\neq\dfrac{25}{8}$. Hence no such factorisation exists and the equation is \textbf{not separable}.
\end{corrige}

\begin{corrige}[Exercise 3 — Definition]
\begin{definition}[First order separable differential equation]
A \emph{first order separable differential equation} is an equation of the form
\[
\frac{dy}{dx}=f(x)\,g(y),
\]
where $f$ is a function of $x$ alone and $g$ is a function of $y$ alone.
\end{definition}

\textbf{Condition for separation.} The variables can be separated precisely when the right-hand side factorises as a \emph{product} (or quotient) of an expression in $x$ only and an expression in $y$ only. In that case, provided $g(y)\neq 0$, one may divide by $g(y)$ and multiply by $dx$ to obtain
\[
\frac{dy}{g(y)}=f(x)\,dx,
\]
in which each side involves only one variable, so that both sides can be integrated independently:
\[
\int \frac{dy}{g(y)}=\int f(x)\,dx + C.
\]
If the right-hand side contains an irreducible sum such as $x+y$ or $\sin(x+y)$, this factorisation is impossible and the equation is not separable.
\end{corrige}

\begin{corrige}[Exercise 4 — Direct integration]
We solve $\dfrac{dy}{dx}=\dfrac{2x}{y}$.

\textbf{Step 1: separate the variables.} Multiplying both sides by $y\,dx$:
\[
y\,dy = 2x\,dx.
\]

\textbf{Step 2: integrate both sides.}
\[
\int y\,dy=\int 2x\,dx
\quad\Longrightarrow\quad
\frac{y^{2}}{2}=x^{2}+C_{1}.
\]

\textbf{Step 3: write the general solution implicitly.} Multiplying by $2$ and setting $C=2C_{1}$:
\[
\boxed{\,y^{2}=2x^{2}+C\,}\qquad (C\in\mathbb{R}).
\]
The integral curves are hyperbolae (for $C\neq 0$) in the $(x,y)$-plane.
\end{corrige}

\begin{corrige}[Exercise 5 — Initial value problem]
We solve $\dfrac{dy}{dx}=\dfrac{3x^{2}}{2y}$, $y(1)=2$.

\textbf{Step 1: separate.}
\[
2y\,dy = 3x^{2}\,dx.
\]

\textbf{Step 2: integrate.}
\[
\int 2y\,dy=\int 3x^{2}\,dx
\quad\Longrightarrow\quad
y^{2}=x^{3}+C.
\]

\textbf{Step 3: apply the initial condition.} With $x=1$, $y=2$:
\[
2^{2}=1^{3}+C \quad\Longrightarrow\quad C=3.
\]

\textbf{Step 4: explicit particular solution.} Since $y(1)=2>0$, we take the positive root:
\[
\boxed{\,y=\sqrt{x^{3}+3}\,}.
\]
\textbf{Check:} $y(1)=\sqrt{1+3}=2$ \checkmark; and $\dfrac{dy}{dx}=\dfrac{3x^{2}}{2\sqrt{x^{3}+3}}=\dfrac{3x^{2}}{2y}$ \checkmark.
\end{corrige}

\begin{corrige}[Exercise 6 — Exponential growth: bacterial culture]
Let $P(t)$ be the population at time $t$ (in hours). Growth at a rate proportional to size gives
\[
\frac{dP}{dt}=kP \qquad (k>0).
\]

\textbf{Step 1: solve.} Separating: $\dfrac{dP}{P}=k\,dt$, hence $\ln P = kt + C$, i.e.
\[
P(t)=P_{0}\,e^{kt},
\]
where $P_{0}=P(0)$.

\textbf{Step 2: determine $k$ from the doubling time.} $P(3)=2P_{0}$:
\[
e^{3k}=2 \quad\Longrightarrow\quad k=\frac{\ln 2}{3}.
\]

\textbf{Step 3: time to triple.} We require $P(t)=3P_{0}$:
\[
e^{kt}=3 \quad\Longrightarrow\quad t=\frac{\ln 3}{k}=\frac{3\ln 3}{\ln 2}.
\]
Numerically,
\[
t=\frac{3\ln 3}{\ln 2}\approx \frac{3\times 1.0986}{0.6931}\approx 4.75\ \text{hours}.
\]
\[
\[\boxed{\,t=\dfrac{3\ln 3}{\ln 2}\approx 4.75\ \text{h}\approx 4\ \text{h}\ 45\ \text{min}.\,}\]
\]
\end{corrige}

\begin{corrige}[Exercise 7 — Mixing problem]
\textbf{(a) Setting up the equation.} Let $S(t)$ be the mass of salt (kg) at time $t$ (min). The volume stays constant at $100$ L (inflow rate $=$ outflow rate $=5$ L/min).

\begin{itemize}
\item \textbf{Salt in:} pure water enters, so rate in $=0$ kg/min.
\item \textbf{Salt out:} the concentration in the tank is $\dfrac{S}{100}$ kg/L, so the rate out is
\[
\frac{S}{100}\times 5=\frac{S}{20}\ \text{kg/min}.
\]
\end{itemize}

Hence
\[
\[\boxed{\;\frac{dS}{dt}=-\frac{S}{20},\qquad S(0)=5.\;}\]
\]

\textbf{(b) Solving.} Separate variables:
\[
\frac{dS}{S}=-\frac{dt}{20}
\quad\Longrightarrow\quad
\ln S=-\frac{t}{20}+C
\quad\Longrightarrow\quad
S=A\,e^{-t/20}.
\]
The initial condition $S(0)=5$ gives $A=5$. Therefore
\[
\[\boxed{\;S(t)=5e^{-t/20}\ \text{kg}.\;}\]
\]
The salt content decays exponentially with time constant $20$ min; after $20$ min only $5e^{-1}\approx 1.84$ kg remain.
\end{corrige}

\begin{corrige}[Exercise 8 — Newton's law of cooling]
We solve $\dfrac{dT}{dt}=-0.07\,(T-20)$, $T(0)=85$.

\textbf{Step 1: separate.} Let $u=T-20$; then $\dfrac{du}{dt}=\dfrac{dT}{dt}$, so
\[
\frac{du}{u}=-0.07\,dt.
\]

\textbf{Step 2: integrate.}
\[
\ln|u|=-0.07t+C \quad\Longrightarrow\quad T-20=Ae^{-0.07t}.
\]

\textbf{Step 3: initial condition.} $T(0)=85$ gives $A=85-20=65$:
\[
T(t)=20+65e^{-0.07t}.
\]

\textbf{Step 4: solve $T=50$.}
\[
50=20+65e^{-0.07t}
\quad\Longrightarrow\quad
e^{-0.07t}=\frac{30}{65}=\frac{6}{13}.
\]
Taking logarithms:
\[
t=\frac{1}{0.07}\ln\!\left(\frac{13}{6}\right)\approx \frac{0.7732}{0.07}\approx 11.0\ \text{min}.
\]
\[
\[\boxed{\,t=\dfrac{100}{7}\ln\!\left(\dfrac{13}{6}\right)\approx 11\ \text{minutes}.\,}\]
\]
\end{corrige}

\begin{corrige}[Exercise 9 — Recognition and substitution (5 marks)]
\textbf{Separability test.} We have
\[
\frac{dy}{dx}=\frac{x^{2}+y^{2}}{xy}=\frac{x}{y}+\frac{y}{x}.
\]
The right-hand side is a sum of terms in $\dfrac{x}{y}$ and $\dfrac{y}{x}$; it cannot be written as a product $f(x)g(y)$ (see the argument in Exercise 2). Hence the equation is \textbf{not separable} as it stands. \hfill[2 marks]

\textbf{Substitution.} The right-hand side depends only on the ratio $\dfrac{y}{x}$:
\[
\frac{x^{2}+y^{2}}{xy}=\frac{1+(y/x)^{2}}{y/x}.
\]
This suggests the \emph{homogeneous substitution}
\[
\[\boxed{\;v=\frac{y}{x},\quad\text{i.e. } y=vx.\;}\]
\]
Then $\dfrac{dy}{dx}=v+x\dfrac{dv}{dx}$, and the equation becomes
\[
v+x\frac{dv}{dx}=\frac{1+v^{2}}{v}
\quad\Longrightarrow\quad
x\frac{dv}{dx}=\frac{1+v^{2}}{v}-v=\frac{1}{v},
\]
that is,
\[
v\,dv=\frac{dx}{x},
\]
which \emph{is} separable, as required. \hfill[3 marks]

\textbf{Examiner's expectations:} clear justification of non-separability (not just a bare assertion); recognition that the equation is homogeneous (function of $y/x$); correct derivative of $y=vx$; explicit display of the separated form.
\end{corrige}

\begin{corrige}[Exercise 10 — General solution and constant solutions (6 marks)]
We solve $\dfrac{dy}{dx}=\dfrac{y\ln y}{x}$ for $x>0$, $y>0$.

\textbf{Step 1: separate.} For $y\ln y\neq 0$,
\[
\frac{dy}{y\ln y}=\frac{dx}{x}. \tag{1 mark}
\]

\textbf{Step 2: integrate.} Using the substitution $u=\ln y$, $du=\dfrac{dy}{y}$:
\[
\int\frac{dy}{y\ln y}=\ln|\ln y|,
\qquad
\int\frac{dx}{x}=\ln x \quad (x>0).
\]
Hence
\[
\ln|\ln y|=\ln x + C_{1}. \tag{2 marks}
\]

\textbf{Step 3: solve for $y$.} Exponentiating, $|\ln y|=e^{C_{1}}x$, so $\ln y = Ax$ where $A\in\mathbb{R}\setminus\{0\}$ (absorbing the sign). Allowing $A=0$ (see below):
\[
\[\boxed{\;y=e^{Ax},\quad A\in\mathbb{R}.\;} \tag{1 mark}\]
\]

\textbf{Step 4: constant solutions lost.} Dividing by $y\ln y$ assumed $y\ln y\neq 0$. Since $y>0$, we need $\ln y\neq 0$, i.e. $y\neq 1$. Check $y\equiv 1$: then $\dfrac{dy}{dx}=0$ and $\dfrac{y\ln y}{x}=\dfrac{1\cdot 0}{x}=0$. So $y=1$ \textbf{is a constant solution}, excluded during separation. \hfill[1 mark]

It is recovered from the general solution by taking $A=0$: $y=e^{0}=1$. Hence the complete family of solutions is $y=e^{Ax}$, $A\in\mathbb{R}$, with $A=0$ giving the constant solution. \hfill[1 mark]

\textbf{Examiner's expectations:} correct separation; recognition of the standard integral $\int\dfrac{dy}{y\ln y}=\ln|\ln y|$; correct handling of the constant of integration; explicit discussion of the solution lost through division.
\end{corrige}

\begin{corrige}[Exercise 11 — Radioactive decay (7 marks)]
\textbf{(a) The model.} Radioactive decay occurs at a rate proportional to the mass present:
\[
\[\boxed{\;\frac{dm}{dt}=-km,\qquad k>0,\quad m(0)=m_{0}.\;} \tag{1 mark}\]
\]

\textbf{(b) Solution.} Separating variables:
\[
\frac{dm}{m}=-k\,dt
\quad\Longrightarrow\quad
\ln m=-kt+C
\quad\Longrightarrow\quad
m=Ae^{-kt}. \tag{1 mark}
\]
Since $m(0)=m_{0}$, $A=m_{0}$, so $m(t)=m_{0}e^{-kt}$.

\textbf{Determination of $k$ from the half-life.} The half-life is $12$ years: $m(12)=\dfrac{m_{0}}{2}$:
\[
e^{-12k}=\frac12
\quad\Longrightarrow\quad
k=\frac{\ln 2}{12}. \tag{2 marks}
\]
Therefore
\[
\[\boxed{\;m(t)=m_{0}\,e^{-\frac{\ln 2}{12}t}=m_{0}\,2^{-t/12}.\;} \tag{1 mark}\]
\]

\textbf{(c) Time to reach 10\% of the original mass.} Set $m(t)=0.1\,m_{0}$:
\[
e^{-kt}=0.1
\quad\Longrightarrow\quad
t=\frac{\ln 10}{k}=\frac{12\ln 10}{\ln 2}. \tag{1 mark}
\]
Numerically,
\[
t=\frac{12\times 2.3026}{0.6931}\approx 39.9\ \text{years}. \tag{1 mark}
\]
\[
\[\boxed{\,t=\dfrac{12\ln 10}{\ln 2}\approx 40\ \text{years}.\,}\]
\]

\textbf{Examiner's expectations:} the negative sign in the model (decay, not growth); correct use of the half-life to fix $k$; exact answer in terms of logarithms plus a sensible decimal approximation; units stated in the final answer.
\end{corrige}

\newpage

\coursec{Summary sheet}

\begin{popfigure}
\begin{tikzpicture}[
  grow cyclic,
  mindmap,
  concept color=popblue!80,
  every node/.style={concept, text=white, font=\small},
  level 1/.append style={level distance=3.4cm, sibling angle=72},
  level 2/.append style={level distance=2.1cm, sibling angle=45, font=\scriptsize}
]
\node[concept, text width=2.6cm, align=center] {\textbf{First Order D.E. with Separable Variables}}
  child[concept color=poporange] { node[concept, text width=2.2cm, align=center] {Meaning of a D.E.}
    child { node[concept, text width=1.9cm, align=center] {order, degree, solution} }
  }
  child[concept color=popgreen!80!black] { node[concept, text width=2.2cm, align=center] {Recognising separable forms}
    child { node[concept, text width=1.9cm, align=center] {$\dfrac{dy}{dx}=f(x)g(y)$} }
  }
  child[concept color=poppurple] { node[concept, text width=2.2cm, align=center] {Separation of variables}
    child { node[concept, text width=1.9cm, align=center] {$\dfrac{dy}{g(y)}=f(x)\,dx$} }
  }
  child[concept color=popgold!90!black] { node[concept, text width=2.2cm, align=center] {General solution}
    child { node[concept, text width=1.9cm, align=center] {integrate both sides $+C$} }
  }
  child[concept color=poppink] { node[concept, text width=2.2cm, align=center] {Particular solution}
    child { node[concept, text width=1.9cm, align=center] {use initial conditions} }
  }
  child[concept color=popdark] { node[concept, text width=2.2cm, align=center] {Applications}
    child { node[concept, text width=1.9cm, align=center] {growth, decay, cooling} }
  };
\end{tikzpicture}
\legende{Mind map of the chapter: first order differential equations with separable variables.}
\end{popfigure}

\begin{bilan}
\textbf{Key definitions.}
\begin{itemize}
  \item A \cle{differential equation} is an equation linking an unknown function $y(x)$, the variable $x$, and derivatives of $y$ such as $\dfrac{dy}{dx}$.
  \item The \cle{order} of a differential equation is the order of the highest derivative present. A \cle{first order} equation contains $\dfrac{dy}{dx}$ but no higher derivative: $F\!\left(x,y,\dfrac{dy}{dx}\right)=0$.
  \item A \cle{solution} is any function $y=\varphi(x)$ which, substituted into the equation, makes it an identity.
  \item The \cle{general solution} contains an arbitrary constant $C$; a \cle{particular solution} is obtained by fixing $C$ using an \cle{initial condition} $y(x_0)=y_0$.
\end{itemize}

\textbf{Recognising separable equations.}
\begin{itemize}
  \item A first order equation is \cle{separable} when it can be written as
  \[ \frac{dy}{dx}=f(x)\,g(y), \]
  i.e.\ the right-hand side is a product of a function of $x$ alone and a function of $y$ alone.
  \item Equivalent forms: $\dfrac{dy}{dx}=\dfrac{f(x)}{h(y)}$, or $M(x)\,dx+N(y)\,dy=0$.
  \item Test: try to factorise the right-hand side into ``$x$-part $\times$ $y$-part''. If impossible (e.g.\ $\dfrac{dy}{dx}=x+y$), the equation is \emph{not} separable.
\end{itemize}

\textbf{Method of separation of variables.}
\begin{enumerate}
  \item Write the equation in the form $\dfrac{dy}{dx}=f(x)g(y)$.
  \item Separate: $\dfrac{dy}{g(y)}=f(x)\,dx$ (all $y$'s with $dy$, all $x$'s with $dx$).
  \item Integrate both sides: $\displaystyle\int\frac{dy}{g(y)}=\int f(x)\,dx$.
  \item Add \emph{one} constant of integration $C$ (one side only is enough).
  \item If possible, make $y$ the subject to obtain the general solution explicitly.
  \item Use the initial condition to find $C$ and state the particular solution.
\end{enumerate}

\textbf{Formulas and standard results to know.}
\begin{itemize}
  \item $\displaystyle\int\frac{dy}{y}=\ln|y|+C$; hence $\dfrac{dy}{dx}=ky \Rightarrow y=Ae^{kx}$ with $A=e^{C}>0$ (or $A\in\mathbb{R}$ allowing sign).
  \item $\displaystyle\int x^{n}\,dx=\dfrac{x^{n+1}}{n+1}+C\ (n\neq-1)$; $\displaystyle\int\frac{dx}{x}=\ln|x|+C$.
  \item Exponential growth/decay law: $\dfrac{dy}{dt}=ky \Rightarrow y=y_0 e^{kt}$ ($k>0$ growth, $k<0$ decay).
  \item Newton's law of cooling: $\dfrac{d\theta}{dt}=-k(\theta-\theta_a)$, giving $\theta-\theta_a=(\theta_0-\theta_a)e^{-kt}$.
  \item When $\ln$ appears, use $e^{\ln|u|+C}=A|u|$ and absorb signs into the arbitrary constant.
\end{itemize}

\textbf{Common mistakes to avoid.}
\begin{itemize}
  \item Forgetting the constant of integration $C$: the general solution is then incomplete and the particular solution cannot be found.
  \item Adding two independent constants $C_1$ and $C_2$; combine them into a single constant $C=C_2-C_1$.
  \item Dividing by $g(y)$ without checking: solutions with $g(y)=0$ (e.g.\ $y=0$) may be lost; mention them when required.
  \item Writing $\ln y$ instead of $\ln|y|$ when the sign of $y$ is unknown.
  \item Separating incorrectly: $\dfrac{dy}{dx}=x+y$ cannot be separated; do not force it.
  \item Algebraic slips when making $y$ the subject, especially with exponentials: $e^{2x+C}=Ae^{2x}$, not $e^{2x}+C$.
  \item Applying the initial condition before integrating, or to the wrong expression.
\end{itemize}

\textbf{GCE examination tips.}
\begin{itemize}
  \item Always show the separation step explicitly: $\dfrac{dy}{g(y)}=f(x)\,dx$; method marks are awarded for it.
  \item State the general solution clearly, then substitute the initial condition in a separate line.
  \item In application problems (population of a town, radioactive decay, cooling of water), define your variables with units, set up the equation $\dfrac{dy}{dt}=ky$ or equivalent, and interpret $k$ from the data.
  \item Give final answers in exact form unless the question asks otherwise; if a numerical value is required, round sensibly (e.g.\ 3 significant figures).
  \item Check your particular solution by substituting it back into the differential equation and the initial condition.
\end{itemize}
\end{bilan}

\findecours

\end{document}
